\documentclass[12pt,a4paper]{ctexbook}
\usepackage{geometry}
\geometry{margin=2.2cm}
\usepackage{setspace}
\setstretch{1.35}
\usepackage{hyperref}
\title{宋词三百首}
\author{古典文献整理}
\date{}
\begin{document}
\maketitle
\tableofcontents
\newpage
\section{湘春夜月}
\textbf{黄孝迈}\par\vspace{0.5em}
近清明。\par
翠禽枝上消魂。\par
可惜一片清歌，都付与黄昏。\par
欲共柳花低诉，怕柳花轻薄，不解伤春。\par
念楚乡旅宿，柔情别绪，谁与温存。\par
空樽夜泣，青山不语，残月当门。\par
翠玉楼前，惟是有、一波湘水，摇荡湘云。\par
天长梦短，问甚时、重见桃根。\par
这次第，算人间没个并刀，翦断心上愁痕。\par

\section{瑞鹤仙}
\textbf{陆叡}\par\vspace{0.5em}
湿云黏雁影。\par
望征路愁迷，离绪难整。\par
千金买光景。\par
但疏钟催晓，乱鸦啼暝。\par
花暗省。\par
许多情、相逢梦境。\par
便行云、都不归来，也合寄将音信。\par
孤迥。\par
盟鸾心在，跨鹤程高，後期无准。\par
情丝待翦。\par
翻惹得，旧时恨。\par
怕天教何处，参差双燕，还染残朱剩粉。\par
对菱花、与说相思，看谁瘦损。\par

\section{渡江云三犯・渡江云}
\textbf{吴文英}\par\vspace{0.5em}
羞红颦浅恨，晚风为落，片绣点重茵。\par
旧堤分燕尾，桂棹轻鸥，宝勒倚残云。\par
千丝怨碧，渐路入、仙坞迷津。\par
肠漫回，隔花时见，背面楚腰身。\par
逡巡。\par
题门惆怅，堕履牵萦，数幽期难准。\par
还始觉、留情缘眼，宽带因春。\par
明朝事与孤烟冷，做满湖、风雨愁人。\par
山黛暝，尘波澹绿无痕。\par

\section{霜叶飞}
\textbf{吴文英}\par\vspace{0.5em}
断烟离绪。\par
关心事，斜阳红隐霜树。\par
半壶秋水荐黄花，香西风雨。\par
纵玉勒、轻飞迅羽。\par
凄凉谁吊荒台古。\par
记醉蹋南屏，彩扇咽、寒蝉倦梦，不知蛮素。\par
聊对旧节传杯，尘笺蠹管，断阕经岁慵赋。\par
小蟾斜影转东篱，夜冷残蛩语。\par
早白发、缘愁万缕。\par
惊飙从卷乌纱去。\par
漫细将、茱萸看，但约明年，翠微高处。\par

\section{瑞鹤仙}
\textbf{吴文英}\par\vspace{0.5em}
晴丝牵绪乱。\par
对沧江斜日，花飞人远。\par
垂杨暗吴苑。\par
正旗亭烟冷，河桥风暖。\par
兰情蕙盼。\par
惹相思、春根酒畔。\par
又争知、吟骨萦销，渐把旧衫重翦。\par
凄断。\par
流红千浪，缺月孤楼，总难留燕。\par
歌尘凝扇。\par
待凭信，拌分钿。\par
试挑灯欲写，还依不忍，笺幅偷和泪卷。\par
寄残云、剩雨蓬莱，也应梦见。\par

\section{宴清都}
\textbf{吴文英}\par\vspace{0.5em}
绣幄鸳鸯柱。\par
红情密，腻云低护秦树。\par
芳根兼倚，花梢钿合，锦屏人妒。\par
东风睡足交枝，正梦花、瑶钗燕股。\par
障滟蜡、满照欢丛，嫠蟾冷落羞度。\par
人间万感幽单，华清惯浴，春盎风露。\par
连鬟并暖，同心共结，向承恩处。\par
凭谁为歌长恨，暗殿锁、秋灯夜语。\par
叙旧期、不负春盟，红朝翠暮。\par

\section{齐天乐}
\textbf{吴文英}\par\vspace{0.5em}
烟波桃叶西陵路，十年断魂潮尾。\par
古柳重攀，轻鸥聚别，陈迹危亭独倚。\par
凉颸乍起。\par
渺烟碛飞帆，暮山横翠。\par
但有江花，共临秋镜照憔悴。\par
华堂烛暗送客，眼波回盼处，芳艳流水。\par
素骨凝冰，柔葱蘸雪，犹忆分瓜深意。\par
清尊未洗。\par
梦不湿行云，漫沾残泪。\par
可惜秋宵，乱蛩疏雨里。\par

\section{风入松}
\textbf{吴文英}\par\vspace{0.5em}
听风听雨过清明。\par
愁草瘗花铭。\par
楼前绿暗分携路，一丝柳、一寸柔情。\par
料峭春寒中酒，交加晓梦啼莺。\par
西园日日扫林亭。\par
依旧赏新晴。\par
黄蜂频扑秋千索，有当时、纤手香凝。\par
惆怅双鸳不到，幽阶一夜苔生。\par

\section{莺啼序}
\textbf{吴文英}\par\vspace{0.5em}
残寒正欺病酒，掩沈香绣户。\par
燕来晚、飞入西城，似说春事迟暮。\par
画船载、清明过却，晴烟冉冉吴宫树。\par
念羁情游荡，随风化为轻絮。\par
十载西湖，傍柳系马，趁娇尘软雾。\par
溯红渐、招入仙溪，锦儿偷寄幽素。\par
倚银屏、春宽梦窄，断红湿、歌纨金缕。\par
暝堤空，轻把斜阳，总还鸥鹭。\par
幽兰旋老，杜若还生，水乡尚寄旅。\par
别後访、六桥无信，事往花委，瘗玉埋香，几番风雨。\par
长波妒盼，遥山羞黛，渔灯分影春江宿，记当时、短楫桃根渡。\par
青楼彷佛，临分败壁题诗，泪墨惨澹尘土。\par
危亭望极，草色天涯，吹鬓侵半苎。\par
暗点检、离痕欢唾，尚染鲛绡，凤迷归，破鸾慵舞。\par
殷勤待写，书中长恨，蓝霞辽海沈过雁，漫相思、弹入哀筝柱。\par
伤心千里江南，怨曲重招，断魂在否。\par

\section{惜黄花慢}
\textbf{吴文英}\par\vspace{0.5em}
送客吴皋。\par
正试霜夜冷，枫落长桥。\par
望天不尽，背城渐杳，离亭黯黯，恨水迢迢。\par
翠香零落红衣老，暮愁锁、残柳眉梢。\par
念瘦腰。\par
沈郎旧日，曾系兰桡。\par
仙人凤咽琼箫。\par
怅断魂送远，九辨难招。\par
醉鬟留盼，小窗翦烛，歌云载恨，飞上银霄。\par
素秋不解随船去，败红趁、一叶寒涛。\par
梦翠翘。\par
怨鸿料过南谯。\par

\section{瑞鹤仙}
\textbf{袁去华}\par\vspace{0.5em}
郊原初过雨。\par
见败叶零乱，风定犹舞。\par
斜阳挂深树。\par
映浓愁浅黛，遥山眉妩。\par
来时旧路。\par
尚岩花、娇黄半吐。\par
到而今，唯有溪边流水，见人如故。\par
无语。\par
邮亭深静，下马还寻，旧曾题处。\par
无聊倦旅。\par
伤离恨，最愁苦。\par
纵收香藏镜，他年重到，人面桃花在否。\par
念沈沈、小阁幽窗，有时梦去。\par

\section{剑器近}
\textbf{袁去华}\par\vspace{0.5em}
夜来雨。\par
赖倩得、东风吹住。\par
海棠正妖饶处。\par
且留取。\par
悄庭户。\par
试细听、莺啼燕语。\par
分明共人愁绪。\par
怕春去。\par
佳树。\par
翠阴初转午。\par
重帘未卷，乍睡起、寂寞看风絮。\par
偷弹清泪寄烟波，见江头故人，为言憔悴如许。\par
彩笺无数。\par
去却寒暄，到了浑无定据。\par
断肠落日千山暮。\par

\section{安公子}
\textbf{袁去华}\par\vspace{0.5em}
弱柳丝千缕。\par
嫩黄匀遍鸦啼处。\par
寒入罗衣春尚浅，过一番风雨。\par
问燕子来时，绿水桥边路。\par
曾画楼、见个人人否。\par
料静掩云窗，尘满哀弦危柱。\par
庾信愁如许。\par
为谁都著眉端聚。\par
独立东风弹泪眼，寄烟波东去。\par
念永昼春闲，人倦如何度。\par
闲傍枕、百啭黄鹂语。\par
唤觉来厌厌，残照依然花坞。\par

\section{卜算子}
\textbf{陆游}\par\vspace{0.5em}
驿外断桥边，寂寞开无主。\par
已是黄昏独自愁，更著风和雨。\par
无意苦争春，一任群芳妒。\par
零落成泥碾作尘，只有香如故。\par

\section{凤箫吟・芳草}
\textbf{韩缜}\par\vspace{0.5em}
锁离愁，连绵无际，来时陌上初熏。\par
绣帏人念远，暗垂珠泪，泣送征轮。\par
长亭长在眼，更重重、远水孤云。\par
但望极楼高，尽日目断王孙。\par
消魂。\par
池塘别後，曾行处、绿妒轻裙。\par
恁时携素手，乱花飞絮里，缓步香。\par
朱颜空自改，向年年、芳意长新。\par
遍绿野，嬉游醉眠，莫负青春。\par

\section{桂枝香}
\textbf{王安石}\par\vspace{0.5em}
登临送目。\par
正故国晚秋，天气初肃。\par
千里澄江似练，翠峰如簇。\par
归帆去棹残阳里，背西风、酒旗斜矗。\par
彩舟云淡，星河鹭起，画图难足。\par
念往昔、繁华竞逐。\par
叹门外楼头，悲恨相续。\par
千古凭高，对此漫嗟荣辱。\par
六朝旧事随流水，但寒烟、芳草凝绿。\par
至今商女，时时犹唱，後庭遗曲。\par

\section{清平乐}
\textbf{王安石}\par\vspace{0.5em}
留春不住。\par
费尽莺儿语。\par
满地残红宫锦污。\par
昨夜南园风雨。\par
小怜初上琵琶。\par
晓来思绕天涯。\par
不肯画堂朱户，春风自在杨花。\par

\section{千秋岁引}
\textbf{王安石}\par\vspace{0.5em}
别馆寒砧，孤城画角。\par
一派秋声入寥廓。\par
东归燕从海上去，南来雁向沙头落。\par
楚台风，庾楼月，宛如昨。\par
无奈被些名利缚。\par
无奈被他情担阁。\par
可惜风流总闲却。\par
当初漫留华表语，而今误我秦楼约。\par
梦阑时，酒醒後，思量著。\par

\section{清平乐}
\textbf{王安国}\par\vspace{0.5em}
留春不住。\par
费尽莺儿语。\par
满地残红宫锦污。\par
昨夜南园风雨。\par
小怜初上琵琶。\par
晓来思绕天涯。\par
不肯画堂朱户，春风自在梨花。\par

\section{临江仙}
\textbf{晏几道}\par\vspace{0.5em}
梦後楼台高锁，酒醒帘幕低垂。\par
去年春恨却来时。\par
落花人独立，微雨燕双飞。\par
记得小苹初见，两重心字罗衣。\par
琵琶弦上说相思。\par
当时明月在，曾照彩云归。\par

\section{蝶恋花}
\textbf{晏几道}\par\vspace{0.5em}
卷絮风头寒欲尽。\par
坠粉飘红，日日香成阵。\par
新酒又添残酒困。\par
今春不减前春恨。\par
蝶去莺飞无处问。\par
隔水高楼，望断双鱼信。\par
恼乱层波横一寸。\par
斜阳只与黄昏近。\par

\section{蝶恋花}
\textbf{晏几道}\par\vspace{0.5em}
醉别西楼醒不记。\par
春梦秋云，聚散真容易。\par
斜月半窗还少睡。\par
画屏闲展吴山翠。\par
衣上酒痕诗里字。\par
点点行行，总是凄凉意。\par
红烛自怜无好计。\par
夜寒空替人垂泪。\par

\section{蝶恋花}
\textbf{晏几道}\par\vspace{0.5em}
梦入江南烟水路。\par
行尽江南，不与离人遇。\par
睡里消魂无说处。\par
觉来惆怅消魂误。\par
欲尽此情书尺素。\par
浮雁沈鱼，终了无凭据。\par
却倚缓弦歌别绪。\par
断肠移破秦筝柱。\par

\section{鹧鸪天}
\textbf{晏几道}\par\vspace{0.5em}
彩袖殷勤捧玉锺。\par
当年拼却醉颜红。\par
舞低杨柳楼心月，歌尽桃花扇影风。\par
从别後，忆相逢。\par
几回魂梦与君同。\par
今宵剩把银釭照，犹恐相逢是梦中。\par

\section{生查子}
\textbf{晏几道}\par\vspace{0.5em}
关山魂梦长，鱼雁音尘少。\par
两鬓可怜青，只为相思老。\par
归梦碧纱窗，说与人人道。\par
真个别离难，不似相逢好。\par

\section{清平乐}
\textbf{晏几道}\par\vspace{0.5em}
留人不住。\par
醉解兰舟去。\par
一棹碧涛春水路。\par
过尽晓莺啼处。\par
渡头杨柳青青。\par
枝枝叶叶离情。\par
此後锦书休寄，画楼云雨无凭。\par

\section{木兰花・玉楼春}
\textbf{晏几道}\par\vspace{0.5em}
秋千院落重帘暮。\par
彩笔闲来题绣户。\par
墙头丹杏雨馀花，门外绿杨风後絮。\par
朝云信断知何处。\par
应作襄王春梦去。\par
紫骝认得旧游踪，嘶过画桥东畔路。\par

\section{菩萨蛮}
\textbf{晏几道}\par\vspace{0.5em}
哀筝一弄湘江曲。\par
声声写尽湘波绿。\par
纤指十三弦。\par
细将幽恨传。\par
当筵秋水慢。\par
玉柱斜飞雁。\par
弹到断肠时。\par
春山眉黛低。\par

\section{玉楼春}
\textbf{晏几道}\par\vspace{0.5em}
东风又作无情计。\par
艳粉娇红吹满地。\par
碧楼帘影不遮愁，还似去年今日意。\par
谁知错管春残事。\par
到处登临曾费泪。\par
此时金盏直须深，看尽落花能几醉。\par

\section{阮郎归}
\textbf{晏几道}\par\vspace{0.5em}
旧香残粉似当初。\par
人情恨不如。\par
一春犹有数行书。\par
秋来书更疏。\par
衾凤冷，枕鸳孤。\par
愁肠待酒舒。\par
梦魂纵有也成虚。\par
那堪和梦无。\par

\section{阮郎归}
\textbf{晏几道}\par\vspace{0.5em}
天边金掌露成霜。\par
云随雁字长。\par
绿杯红袖称重阳。\par
人情似故乡。\par
兰佩紫，菊簪黄。\par
殷勤理旧狂。\par
欲将沈醉换悲凉。\par
清歌莫断肠。\par

\section{六么令}
\textbf{晏几道}\par\vspace{0.5em}
绿阴春尽，飞絮绕香阁。\par
晚来翠眉宫样，巧把远山学。\par
一寸狂心未说，已向横波觉。\par
画帘遮币。\par
新翻曲妙，暗许闲人带偷掏。\par
前度书多隐语，意浅愁难答。\par
昨夜诗有回纹，韵险还慵押。\par
都待笙歌散了，记取留时霎。\par
不消红蜡。\par
闲云归後，月在庭花旧阑角。\par

\section{御街行}
\textbf{晏几道}\par\vspace{0.5em}
街南绿树春饶絮。\par
雪满游春路。\par
树头花艳杂娇云，树底人家朱户。\par
北楼闲上，疏帘高卷，直见街南树。\par
阑干倚尽犹慵去。\par
几度黄昏雨。\par
晚春盘马踏青苔，曾傍绿阴深驻。\par
落花犹在，香屏空掩，人面知何处。\par

\section{虞美人}
\textbf{晏几道}\par\vspace{0.5em}
曲阑干外天如水。\par
昨夜还曾倚。\par
初将明月比佳期。\par
长向月圆时候、望人归。\par
罗衣著破前香在。\par
旧意谁教改。\par
一春离恨懒调弦。\par
犹有两行闲泪、宝筝前。\par

\section{留春令}
\textbf{晏几道}\par\vspace{0.5em}
画屏天畔，梦回依约，十洲云水。\par
手拈红笺寄人书，写无限、伤春事。\par
别浦高楼曾漫倚。\par
对江南千里。\par
楼下分流水声中，有当日、凭高泪。\par

\section{思远人}
\textbf{晏几道}\par\vspace{0.5em}
红叶黄花秋意晚，千里念行客。\par
飞云过尽，归鸿无信，何处寄书得。\par
泪弹不尽临窗滴。\par
就砚旋研墨。\par
渐写到别来，此情深处，红笺为无色。\par

\section{生查子}
\textbf{王观}\par\vspace{0.5em}
关山魂梦长，塞雁音书少。\par
两鬓可怜青，一夜相思老。\par
归傍碧纱窗，说与人人道。\par
真个别离难，不似相逢好。\par

\section{水龙吟}
\textbf{苏轼}\par\vspace{0.5em}
似花还似非花，也无人惜从教坠。\par
抛家傍路，思量却是，无情有思。\par
萦损柔肠，困酣娇眼，欲开还闭。\par
梦随风万里，寻郎去处，又还被、莺呼起。\par
不恨此花飞尽，恨西园、落红难缀。\par
晓来雨过，遗踪何在，一池萍碎。\par
春色三分，二分尘土，一分流水。\par
细看来，不是杨花点点，是离人泪。\par

\section{水调歌头}
\textbf{苏轼}\par\vspace{0.5em}
明月几时有，把酒问青天。\par
不知天上宫阙，今夕是何年。\par
我欲乘风归去，又恐琼楼玉宇，高处不胜寒。\par
起舞弄清影，何似在人间。\par
转朱阁，低绮户，照无眠。\par
不应有恨，何事长向别时圆。\par
人有悲欢离合，月有阴晴圆缺，此事古难全。\par
但愿人长久，千里共婵娟。\par

\section{临江仙}
\textbf{苏轼}\par\vspace{0.5em}
夜饮东坡醒复醉，归来仿佛三更。\par
家童鼻息已雷鸣。\par
敲门都不应，倚杖听江声。\par
长恨此身非我有，何时忘却营营。\par
夜阑风静縠纹平。\par
小舟从此逝，江海寄馀生。\par

\section{定风波}
\textbf{苏轼}\par\vspace{0.5em}
莫听穿林打叶声。\par
何妨吟啸且徐行。\par
竹杖芒鞋轻胜马。\par
谁怕。\par
一蓑烟雨任平生。\par
料峭春风吹酒醒。\par
微冷。\par
山头斜照却相迎。\par
回首向来萧瑟处。\par
归去。\par
也无风雨也无晴。\par

\section{卜算子}
\textbf{苏轼}\par\vspace{0.5em}
缺月挂疏桐，漏断人初静。\par
时见幽人独往来，缥缈孤鸿影。\par
惊起却回头，有恨无人省。\par
拣尽寒枝不肯栖，枫落吴江冷。\par

\section{贺新郎}
\textbf{苏轼}\par\vspace{0.5em}
乳燕飞华屋。\par
悄无人、桐阴转午，晚凉新浴。\par
手弄生绡白团扇，扇手一时似玉。\par
渐困倚、孤眠清熟。\par
帘外谁来推绣户，枉教人、梦断瑶台曲。\par
又却是，风敲竹。\par
石榴半吐红巾蹙。\par
待浮花、浪蕊都尽，伴君幽独。\par
艳一枝细看取，芳心千重似束。\par
又恐被、秋风惊绿。\par
若待得君来向此，花前对酒不忍触。\par
共粉泪，两蔌蔌。\par

\section{洞仙歌}
\textbf{苏轼}\par\vspace{0.5em}
冰肌玉骨，自清凉无汗。\par
水殿风来暗香满。\par
绣帘开、一点明月窥人，人未寝、枕钗横鬓乱。\par
起来携素手，庭户无声，时见疏星渡河汉。\par
试问夜如何，夜已三更，金波淡、玉绳低转。\par
但屈指、西风几时来，又不道、流年暗中偷换。\par

\section{江神子・江城子}
\textbf{苏轼}\par\vspace{0.5em}
十年生死两茫茫。\par
不思量。\par
自难忘。\par
千里孤坟，无处话凄凉。\par
纵使相逢应不识，尘满面，鬓如霜。\par
夜来幽梦忽还乡。\par
小轩窗。\par
正梳妆。\par
相顾无言，惟有泪千行。\par
料得年年断肠处，明月夜，短松冈。\par

\section{永遇乐}
\textbf{苏轼}\par\vspace{0.5em}
明月如霜，好风如水，清景无限。\par
曲港跳鱼，圆荷泻露，寂寞无人见。\par
如三鼓，铿然一叶，黯黯梦云惊断。\par
夜茫茫，重寻无处，觉来小园行遍。\par
天涯倦客，山中归路，望断故园心眼。\par
燕子楼空，佳人何在，空锁楼中燕。\par
古今如梦，何曾梦觉，但有旧欢新怨。\par
异时对，黄楼夜景，为余浩叹。\par

\section{青玉案}
\textbf{苏轼}\par\vspace{0.5em}
三年枕上吴中路。\par
遣黄耳、随君去。\par
若到松江呼小渡。\par
莫惊鸥鹭，四桥尽是，老子经行处。\par
辋川图上看春暮。\par
常记高人右丞句。\par
作个归期天已许。\par
春衫犹是，小蛮针线，曾湿西湖雨。\par

\section{谢池春}
\textbf{李之仪}\par\vspace{0.5em}
残寒销尽，疏雨过、清明後。\par
花径敛馀红，风沼萦新皱。\par
乳燕穿庭户，飞絮沾襟袖。\par
正佳时，仍晚昼。\par
著人滋味，真个浓如酒。\par
频移带眼，空只恁、厌厌瘦。\par
不见又思量，见了还依旧。\par
为问频相见，何似长相守。\par
天不老，人未偶。\par
且将此恨，分付庭前柳。\par

\section{卜算子}
\textbf{李之仪}\par\vspace{0.5em}
我住长江头，君住长江尾。\par
日日思君不见君，共饮长江水。\par
此水几时休，此恨何时已。\par
只愿君心似我心，定不负相思意。\par

\section{虞美人}
\textbf{舒亶}\par\vspace{0.5em}
芙蓉落尽天涵水。\par
日暮沧波起。\par
背飞双燕帖云寒。\par
独向小楼东畔、倚阑看。\par
浮生只合尊前老。\par
雪满长安道。\par
故人早晚上高台。\par
赠我江南春色、一枝梅。\par

\section{高阳台}
\textbf{韩}\par\vspace{0.5em}
频听银签，重燃绛蜡，年华衮衮惊心。\par
饯旧迎新，能消几刻光阴。\par
老来可惯通宵饮，待不眠、还怕寒侵。\par
掩清尊。\par
多谢梅花，伴我微吟。\par
邻娃已试春妆了，更蜂腰簇翠，燕股横金。\par
勾引东风，也知芳思难禁。\par
朱颜那有年年好，逞艳游、赢取如今。\par
恣登临。\par
残雪楼台，迟日园林。\par

\section{玉楼春}
\textbf{严仁}\par\vspace{0.5em}
春风只在园西畔。\par
荠菜花繁胡蝶乱。\par
冰池晴绿照还空，香径落红吹已断。\par
意长翻恨游丝短。\par
尽日相思罗带缓。\par
宝奁明月不欺人，明日归来君试看。\par

\section{生查子}
\textbf{刘克庄}\par\vspace{0.5em}
繁灯夺霁华。\par
戏鼓侵明发。\par
物色旧时同，情味中年别。\par
浅画镜中眉，深拜楼西月。\par
人散市声收，渐入愁时节。\par

\section{高阳台}
\textbf{张炎}\par\vspace{0.5em}
接叶巢莺，平波卷絮，断桥斜日归船。\par
能几番游，看花又是明年。\par
东风且伴蔷薇住，到蔷薇、春已堪怜。\par
更凄然。\par
万绿西泠，一抹荒烟。\par
当年燕子知何处，但苔深韦曲，草暗斜川。\par
见说新愁，如今也到鸥边。\par
无心再续笙歌梦，掩重门、浅醉闲眠。\par
莫开帘。\par
怕见飞花，怕听啼鹃。\par

\section{贺新郎}
\textbf{刘克庄}\par\vspace{0.5em}
深院榴花吐。\par
画帘开、衣纨扇，午风清暑。\par
儿女纷纷夸结束，新样钗符艾虎。\par
早已有、游人观渡。\par
老大逢场慵作戏，任陌头、年少争旗鼓。\par
溪雨急，浪花舞。\par
灵均标致高如许。\par
忆生平、既纫兰佩，更怀椒糈。\par
谁信骚魂千载後，波底垂涎角黍。\par
又说是、蛟馋龙怒。\par
把似而今醒到了，料当年、醉死差无苦。\par
聊一笑，吊千古。\par

\section{贺新郎}
\textbf{刘克庄}\par\vspace{0.5em}
湛湛长空黑。\par
更那堪、斜风细雨，乱愁如织。\par
老眼平生空四海，赖有高楼百尺。\par
看浩荡、千崖秋色。\par
白发书生神州泪，尽凄凉、不向牛山滴。\par
追往事，去无迹。\par
少年自负凌云笔。\par
到而今、春华落尽，满怀萧瑟。\par
常恨世人新意少，爱说南朝狂客。\par
把破帽、年年拈出。\par
若对黄花孤负酒，怕黄花、也笑人岑寂。\par
鸿北去，日西匿。\par

\section{玉楼春}
\textbf{刘克庄}\par\vspace{0.5em}
年年跃马长安市。\par
客舍似家家似寄。\par
青钱换酒日无何，红烛呼庐宵不寐。\par
易挑锦妇机中字。\par
难得玉人心下事。\par
男儿西北有神州，莫滴水西桥畔泪。\par

\section{好事近}
\textbf{韩元吉}\par\vspace{0.5em}
凝碧旧池头，一听管弦凄切。\par
多少梨园声在，总不堪华发。\par
杏花无处避春愁，也傍野烟发。\par
惟有御沟声断，似知人呜咽。\par

\section{六州歌头}
\textbf{韩元吉}\par\vspace{0.5em}
东风著意。\par
先上小桃枝。\par
红粉腻。\par
娇如醉。\par
倚朱扉。\par
记年时。\par
隐映新妆，面临水岸。\par
春将半。\par
云日暖。\par
斜桥转。\par
夹城西。\par
草软莎平跋马，垂杨渡、玉勒争嘶。\par
认蛾眉凝笑，脸薄拂燕支。\par
绣户曾窥。\par
恨依依。\par
共携手处。\par
香如雾。\par
红随步。\par
怨春迟。\par
消瘦损。\par
凭谁问。\par
只花知。\par
泪空垂。\par
旧日堂前燕，和烟雨，又双飞。\par
人自老。\par
春长好。\par
梦佳期。\par
前度刘郎，几许风流地，花也应悲。\par
但茫茫暮霭，目断武陵溪。\par
往事难追。\par

\section{烛影摇红}
\textbf{张抡}\par\vspace{0.5em}
双阙中天，凤楼十二春寒浅。\par
去年元夜奉宸游，曾侍瑶池宴。\par
玉殿珠帘尽卷。\par
拥群仙、蓬壶阆苑。\par
五云深处，万烛光中，揭天丝管。\par
驰隙流年，恍如一瞬星霜换。\par
今宵谁念泣孤臣，回首长安远。\par
可是尘缘未断。\par
谩惆怅、华胥梦短。\par
满怀幽恨，数点寒灯，几声归雁。\par

\section{永遇乐}
\textbf{辛弃疾}\par\vspace{0.5em}
千古江山，英雄无觅，孙仲谋处。\par
舞榭歌台，风流总被，雨打风吹去。\par
斜阳草树，寻常巷陌，人道寄奴曾住。\par
想当年，金戈铁马，气吞万里如虎。\par
元嘉草草，封狼居胥，赢得仓皇北顾。\par
四十三年，望中犹记，烽火扬州路。\par
可堪回首，佛狸祠下，一片神鸦社鼓。\par
凭谁问，廉颇老矣，尚能饭否。\par

\section{绿意・疏影}
\textbf{张炎}\par\vspace{0.5em}
碧圆自洁。\par
向浅洲远渚，亭亭清绝。\par
犹有遗簪，不展秋心，能卷几多炎热。\par
鸳鸯密语同倾盖，且莫与、浣纱人说。\par
恐怨歌、忽断花风，碎却翠云千叠。\par
回首当年汉舞，怕飞去、谩皱留仙裙摺。\par
恋恋青衫，犹染枯香，还叹鬓丝飘雪。\par
盘心清露如铅水，又一夜、西风吹折。\par
喜静看、匹练秋光，倒泻半湖明月。\par

\section{绿头鸭・多丽}
\textbf{晁端礼}\par\vspace{0.5em}
晚云收，淡天一片琉璃。\par
烂银盘、来从海底，皓色千里澄辉。\par
莹无尘、素娥淡伫，静可数、丹桂参差。\par
玉露初零，金风未凛，一年无似此佳时。\par
露坐久，疏莹时度，乌鹊正南飞。\par
瑶台冷，栏干凭暖，玉下迟迟。\par
念佳人，音尘别後，对此应解相思。\par
最关情、漏声正永，暗断肠、花影偷移。\par
料得来宵，清光未减，阴晴天气又争知。\par
共凝恋、如今别後，还是隔年期。\par
人强健，清尊素影，长愿相随。\par

\section{洞仙歌}
\textbf{李元膺}\par\vspace{0.5em}
雪云散尽，放晓晴池院。\par
杨柳于人便青眼。\par
更风流多处，一点梅心、相映远。\par
约略轻笑浅。\par
一年春好处，不在浓芳，小艳疏香最娇软。\par
到清明时候，百紫千红花正乱。\par
已失春风一半。\par
蚤占取韶光，共追游，但莫管春寒，醉红自暖。\par

\section{渔家傲}
\textbf{朱服}\par\vspace{0.5em}
小雨廉纤风细细。\par
万家杨柳轻烟里。\par
恋树湿花飞不起。\par
愁无比。\par
和春付与西流水。\par
九十光阴能有几。\par
金龟解尽留无计。\par
寄语东城沽酒市。\par
拼一醉。\par
而今乐事他年泪。\par

\section{青门饮}
\textbf{时彦}\par\vspace{0.5em}
胡马嘶风，汉旗翻雪，彤云又吐，一竿残照。\par
古木连空，乱山无数，行尽暮沙衰草。\par
星斗横幽馆，夜无眠、灯花空老。\par
雾浓香鸭，冰凝泪烛，霜天难晓。\par
长记晓妆才了，一杯未尽，离怀多少。\par
醉里秋波，梦中朝雨，都是醒时烦恼。\par
料有牵情处，忍思量、耳边曾道。\par
甚时跃马归来，认得迎门轻笑。\par

\section{望海潮}
\textbf{秦观}\par\vspace{0.5em}
梅英疏淡，冰澌溶泄，东风暗换年华。\par
金谷俊游，铜驼巷陌，新晴细履平沙。\par
长记误随车。\par
正絮翻蝶舞，芳思交加。\par
柳下桃蹊，乱分春色到人家。\par
西园夜饮鸣笳。\par
有华灯碍月，飞盖妨花。\par
兰苑未空，行人渐老，重来是事堪嗟。\par
烟暝酒旗斜。\par
但倚楼极目，时见栖鸦。\par
无奈归心。\par
暗随流水到天涯。\par

\section{满庭芳}
\textbf{秦观}\par\vspace{0.5em}
山抹微云，天连衰草，画角声断谯门。\par
暂停征棹，聊共引离尊。\par
多少蓬莱旧事，空回首、烟霭纷纷。\par
斜阳外，寒鸦万点，流水绕孤村。\par
销魂。\par
当此际，香囊暗解，罗带轻分。\par
谩赢得、青楼薄名存。\par
此去何时见也，襟袖上、空惹啼痕。\par
伤情处，高城望断，灯火已黄昏。\par

\section{浣溪沙}
\textbf{秦观}\par\vspace{0.5em}
漠漠轻寒上小楼。\par
晓阴无赖似穷秋。\par
淡烟流水画屏幽。\par
自在飞花轻似梦，无边丝雨细如愁。\par
宝帘闲挂小银钩。\par

\section{阮郎归}
\textbf{秦观}\par\vspace{0.5em}
湘天风雨破寒初。\par
深沈庭院虚。\par
丽谯吹罢《小单于》。\par
迢迢清夜徂。\par
乡梦断，旅魂孤。\par
峥嵘岁又除。\par
衡阳犹有雁传书。\par
郴阳和雁无。\par

\section{满庭芳}
\textbf{秦观}\par\vspace{0.5em}
晓色云开，春随人意，骤雨才过还晴。\par
古台芳榭，飞燕蹴红英。\par
舞困榆钱自落，秋千外、绿水桥平。\par
东风里，朱门映柳，低底按小秦筝。\par
多情。\par
行乐处，珠钿翠盖，玉辔红缨。\par
渐酒空金，花困蓬瀛。\par
豆蔻梢头旧恨，十年梦，屈指堪惊。\par
凭阑久，疏烟淡日，寂寞下芜城。\par

\section{帝台春}
\textbf{李甲}\par\vspace{0.5em}
芳草碧色，萋萋遍南陌。\par
暖絮乱红，也知人、春愁无力。\par
忆得盈盈拾翠侣，共携赏、凤城寒食。\par
到今来，海角逢春，天涯为客。\par
愁旋释。\par
还似织。\par
泪暗拭。\par
又偷滴。\par
谩伫立、遍倚危阑，尽黄昏，也只是、暮云凝碧。\par
拼则而已今拼了，忘则怎生便忘得。\par
又还问鳞鸿，试重寻消息。\par

\section{蝶恋花}
\textbf{赵令}\par\vspace{0.5em}
欲减罗衣寒未去。\par
不卷珠帘，人在深深处。\par
红杏枝头花几许。\par
啼痕止恨清明雨。\par
尽日沈烟香一缕。\par
宿酒醒迟，恼破春情绪。\par
飞燕又将归信误。\par
小屏风上西江路。\par

\section{醉中真・浣溪沙}
\textbf{贺铸}\par\vspace{0.5em}
不信芳春厌老人。\par
老人几度送春归。\par
惜春行乐莫辞频。\par
巧笑艳歌皆我意，恼花颠酒拼君嗔。\par
物情惟有醉中真。\par

\section{人南渡・感皇恩}
\textbf{贺铸}\par\vspace{0.5em}
兰芷满芳洲，游思横路。\par
罗袜尘生步。\par
迎顾。\par
整鬟颦黛，脉脉两情难语。\par
细风吹柳絮。\par
人南渡。\par
回首旧游，山无重数。\par
花底深朱户。\par
何处。\par
半黄梅子，向晚一帘疏雨。\par
断魂分付与。\par
春将去。\par

\section{薄幸}
\textbf{贺铸}\par\vspace{0.5em}
艳真多态。\par
更的的，频回眄睐。\par
便认得、琴心相许，与写宜男双带。\par
记画堂、斜月朦胧，轻颦微笑娇无奈。\par
便翡翠屏开，芙蓉帐掩，与把香罗偷解。\par
自过了收灯後，都不见、踏青挑菜。\par
几回凭双燕，丁宁深意，往来翻恨重帘碍。\par
约何时再。\par
正春浓酒暖，人闲昼永无聊赖。\par
厌厌睡起，犹有花梢日在。\par

\section{伴云来・天香}
\textbf{贺铸}\par\vspace{0.5em}
烟络横林，山沈远照，逦迤黄昏钟鼓。\par
烛映帘栊，蛩催机杼，共苦清秋风露。\par
不眠思妇，齐应和、几声砧杵。\par
惊动天涯倦宦，岁华行暮。\par
当年酒狂自负。\par
谓东君、已春相付。\par
流浪征骖北道，客墙南浦。\par
幽恨无人晤语。\par
赖明月曾知旧游处。\par
好伴云来，还将梦去。\par

\section{青玉案}
\textbf{无名氏}\par\vspace{0.5em}
一年春事都来几。\par
早过了、三之二。\par
绿暗红嫣浑可事。\par
绿杨庭院，暖风帘幕，有个人憔悴。\par
买花载酒长安市。\par
又争似家山见桃李。\par
不枉东风吹客泪，相思难表，梦魂无据，惟有归来是。\par

\section{高阳台}
\textbf{吴文英}\par\vspace{0.5em}
修竹凝妆，垂杨驻马，凭阑浅画成图。\par
山色谁题，楼前有雁斜书。\par
东风紧送斜阳下，弄旧寒、晚酒醒馀。\par
自销凝，能几花前，顿老相如。\par
伤春不在高楼上，在灯前攲枕，雨外熏炉。\par
怕舣游船，临流可奈清臞。\par
飞红若到西湖底，搅翠澜、总是愁鱼。\par
莫重来，吹尽香绵，泪满平芜。\par

\section{高阳台}
\textbf{吴文英}\par\vspace{0.5em}
宫粉雕痕，仙云堕影，无人野水荒湾。\par
古石埋香，金沙锁骨连环。\par
南楼不恨吹横笛，恨晓风、千里关山。\par
半飘零，庭上黄昏，月冷阑干。\par
寿阳空理愁鸾。\par
问谁调玉髓，暗补香瘢。\par
细雨归鸿，孤山无限春寒。\par
离魂难倩招清些，梦缟衣、解佩溪边。\par
最愁人，啼鸟晴明，叶底青圆。\par

\section{三姝媚}
\textbf{吴文英}\par\vspace{0.5em}
湖山经醉惯。\par
溃春衫、啼痕酒痕无限。\par
又客长安，叹断襟零袂，尘谁浣。\par
紫曲门荒，沿败井、风摇青蔓。\par
对语东邻，犹是曾巢，谢堂双燕。\par
春梦人间须断。\par
但怪得、当年梦缘能短。\par
绣屋秦筝，傍海棠偏爱，夜深开宴。\par
舞歇歌沈，花未减、红颜先变。\par
伫久河桥欲去，斜阳泪满。\par

\section{八声甘州}
\textbf{吴文英}\par\vspace{0.5em}
渺空烟四远，是何年、青天坠长星。\par
幻苍崖云树，名娃金屋，残霸宫城。\par
箭径酸风射眼，腻水染花腥。\par
时靸双鸳响，廊叶秋声。\par
宫里吴王沈醉，倩五湖倦客，独钓醒醒。\par
问苍波无语，华发奈山青。\par
水涵空、阑干高处，送乱鸦、斜日落渔汀。\par
连呼酒，上琴台去，秋与云平。\par

\section{夜合花}
\textbf{吴文英}\par\vspace{0.5em}
柳暝河桥，莺晴台苑，短策频惹春香。\par
当时夜泊，温柔便入深乡。\par
词韵窄，酒杯长。\par
翦蜡花、壶箭催忙。\par
共追游处，凌波翠陌，连棹横塘。\par
十年一梦凄凉。\par
似西湖燕去，吴馆巢荒。\par
重来万感，依前唤酒银罂。\par
溪雨急，岸花狂。\par
趁残鸦、飞过苍茫。\par
故人楼上，凭谁指与，芳草斜阳。\par

\section{蹋莎行・踏莎行}
\textbf{吴文英}\par\vspace{0.5em}
润玉笼绡，檀樱倚扇。\par
绣圈犹带脂香浅。\par
榴心空叠舞裙红，艾枝应压愁鬟乱。\par
午梦千山，窗阴一箭。\par
香瘢新褪红丝腕。\par
隔江人在雨声中，晚风菰叶生秋怨。\par

\section{夜游宫}
\textbf{吴文英}\par\vspace{0.5em}
人去西楼雁杳。\par
叙别梦、扬州一觉。\par
云澹星疏楚山晓。\par
听啼乌，立河桥，话未了。\par
雨外蛩声早。\par
细织就、霜丝多少。\par
说与萧娘未知道。\par
向长安，对秋灯，几人老。\par

\section{鹧鸪天}
\textbf{吴文英}\par\vspace{0.5em}
池上红衣伴倚阑。\par
栖鸦常带夕阳还。\par
殷云度雨疏桐落，明月生凉宝扇闲。\par
乡梦窄，水天宽。\par
小窗愁黛澹秋山。\par
吴鸿好为传归信，杨柳阊门屋数间。\par

\section{唐多令}
\textbf{吴文英}\par\vspace{0.5em}
何处合成愁。\par
离人心上秋。\par
纵芭蕉、不雨也飕飕。\par
都道晚凉天气好，有明月、怕登楼。\par
年事梦中休。\par
花空烟水流。\par
燕辞归、客尚淹留。\par
垂柳不萦裙带住，漫长是、系行舟。\par

\section{金缕歌・贺新郎}
\textbf{吴文英}\par\vspace{0.5em}
乔木生云气。\par
访中兴、英雄陈迹，暗追前事。\par
战舰东风悭借便，梦断神州故里。\par
旋小筑、吴宫闲地。\par
华表月明归夜鹤，叹当时、花竹今如此。\par
枝上露，溅清泪。\par
遨头小簇行春队。\par
步苍苔、寻幽别坞，问梅开未。\par
重唱梅边吸度曲，催发寒梢冻蕊。\par
此心与、东君同意。\par
後不如今非昔，两无言、相对沧浪水。\par
怀此恨，寄残醉。\par

\section{大有}
\textbf{潘希白}\par\vspace{0.5em}
戏马台前，采花篱下，问岁华、还是重九。\par
恰归来、南山翠色依旧。\par
帘栊昨夜听风雨，都不似、登归时候。\par
一片宋玉情怀，十分卫郎清瘦。\par
红萸佩、空对酒。\par
砧杆动微寒，暗欺罗袖。\par
秋已无多，早是败荷衰柳。\par
强整帽檐侧，曾经向、天涯搔首。\par
几回忆、故国莼鲈，霜前雁後。\par

\section{兰陵王}
\textbf{刘辰翁}\par\vspace{0.5em}
送春去。\par
春去人间无路。\par
秋千外、芳草连天，谁遣风沙暗南浦。\par
依依甚意绪。\par
漫忆海门飞絮。\par
乱鸦过，斗转城荒，不见来时试灯处。\par
春去。\par
最谁苦。\par
但箭雁沉边，梁燕无主。\par
杜鹃声里长门暮。\par
想玉树凋土，泪盘如露。\par
咸阳送客屡回顾。\par
斜日未能度。\par
春去。\par
尚来否。\par
正江令恨别，庾信愁赋。\par
苏堤尽日风和雨。\par
叹神游故国，花记前度。\par
人生流落，顾孺子，共夜雨。\par

\section{宝鼎现}
\textbf{刘辰翁}\par\vspace{0.5em}
红妆春骑。\par
踏月影、竿旗穿市。\par
望不尽、楼台歌舞，习习香尘莲步底。\par
箫声断、约彩鸾归去，未怕金吾呵醉。\par
甚辇路、喧阗且止。\par
听得念奴歌起。\par
父老犹记宣和事。\par
抱铜仙、清泪如水。\par
还转盼、沙河多丽。\par
漾明光连邸第。\par
帘影冻、散红光成绮。\par
月浸葡萄十里。\par
看往来、神仙才子。\par
肯把菱花扑碎。\par
肠断竹马儿童，空见说、三千乐指。\par
等多时春不归来，到春时欲睡。\par
又说向、灯前拥髻。\par
暗滴鲛珠坠。\par
便当日、亲见霓裳，天上人间梦里。\par

\section{永遇乐}
\textbf{刘辰翁}\par\vspace{0.5em}
璧月初晴，黛云远澹，春事谁主。\par
禁苑娇寒，湖堤倦暖，前度遽如许。\par
香尘暗陌，华灯明昼，长是懒携手去。\par
谁知道，断烟禁夜，满城似愁风雨。\par
宣和旧日，临安南渡，芳景犹自如故。\par
缃帙流离，风鬟三五，能赋词最苦。\par
江南无路，州今夜，此苦又谁知否。\par
空相对，残无寐，满村社鼓。\par

\section{摸鱼儿}
\textbf{刘辰翁}\par\vspace{0.5em}
怎知他、春归何处，相逢且尽尊酒。\par
少年袅袅天涯恨，长结西湖烟柳。\par
休回首。\par
但细雨断桥，憔悴人归後。\par
东风似旧。\par
问前度桃花，刘郎能记，花复认郎否。\par
君且住，草草留君翦韭。\par
前宵更恁时候。\par
深杯欲共歌声滑，翻湿春衫半袖。\par
空眉皱。\par
看白发尊前，已似人人有。\par
临分把手。\par
叹一笑论文，清狂顾曲，此会几时又。\par

\section{绣鸾凤花犯・花犯}
\textbf{周密}\par\vspace{0.5em}
楚江湄，湘娥乍见，无言洒清泪。\par
淡然春意。\par
空独倚东风，芳思谁寄。\par
凌波路冷秋无际。\par
香云随步起。\par
谩记得，汉宫仙掌，亭亭明月底。\par
冰弦写怨更多情，骚人恨，枉赋芳兰幽芷。\par
春思远，谁叹赏、国香风味。\par
相将共、岁寒伴侣。\par
小窗净、沈烟熏翠袂。\par
幽梦觉，涓涓清露，一枝灯影里。\par

\section{瑶花慢}
\textbf{周密}\par\vspace{0.5em}
朱钿宝。\par
天上飞琼，比人间春别。\par
江南江北，曾未见，谩拟梨云梅雪。\par
淮山春晚，问谁识、芳心高洁。\par
消几番、花落花开，老了玉关豪杰。\par
金壶翦送琼枝，看一骑红尘，香度瑶阙。\par
韶华正好，应自喜、初识长安蜂蝶。\par
杜郎老矣，想旧事、花须能说。\par
记少年，一梦扬州，二十四桥明月。\par

\section{玉京秋}
\textbf{周密}\par\vspace{0.5em}
烟水阔。\par
高林弄残照，晚蜩凄切。\par
碧砧度韵，银床飘叶。\par
衣湿桐阴露冷，采凉花、时赋秋雪。\par
叹轻别。\par
一襟幽事，砌蛩能说。\par
客思吟商还怯。\par
怨歌长、琼壶暗缺。\par
翠扇恩疏，红衣香褪，翻成消歇。\par
玉骨西风，恨最恨、闲却新凉时节。\par
楚箫咽。\par
谁倚西楼淡月。\par

\section{曲游春}
\textbf{周密}\par\vspace{0.5em}
楚苑东风外，暖丝情絮，春思如织。\par
燕约莺期，恼芳情偏在，翠深红隙。\par
漠漠香尘隔。\par
沸十里、乱弦丛笛。\par
看画船，尽入西泠，闲却半湖春色。\par
柳陌。\par
新烟凝碧。\par
映帘底宫眉，堤上游勒。\par
轻暝笼寒，怕梨云梦冷，杏香愁幂。\par
歌管酬寒食。\par
奈蝶怨、良宵岑寂。\par
正满湖、碎月摇花，怎生去得。\par

\section{高阳台}
\textbf{周密}\par\vspace{0.5em}
照野旌旗，朝天车马，平沙万里天低。\par
宝带金章，尊前茸帽风。\par
秦关汴水经行地，想登临、都付新诗。\par
纵英游，叠鼓清笳，骏马名姬。\par
酒酣应对燕山雪，正冰河月冻，晓陇云飞。\par
投老残年，江南谁念方回。\par
东风渐绿西湖柳，雁已还、人未南归。\par
最关情，折尽梅花，难寄相思。\par

\section{摸鱼儿}
\textbf{朱嗣发}\par\vspace{0.5em}
对西风、鬓摇烟碧，参差前事流水。\par
紫丝罗带鸳鸯结，的的镜盟钗誓。\par
浑不记、漫手织回文，几度欲心碎。\par
安花著蒂。\par
奈雨覆云翻，情宽分窄，石上玉簪脆。\par
朱楼外。\par
愁压空云欲坠。\par
月痕犹照无寐。\par
阴晴也只随天意。\par
枉了玉消香碎。\par
君且醉。\par
君不见、长门青草春风泪。\par
一时左计。\par
悔不早荆钗，暮天修竹，头白倚寒翠。\par

\section{解佩环・疏影}
\textbf{彭元逊}\par\vspace{0.5em}
江空不渡，恨蘼芜杜若，零落无数。\par
远道荒寒，婉娩流年，望望美人迟暮。\par
风烟雨雪阴晴晚，更何须，春风千树。\par
尽孤城、落木萧萧，日夜江声流去。\par
日晏山深闻笛，恐他年流落，与子同赋。\par
事阔心违，交淡媒劳，蔓草沾衣多露。\par
汀洲窈窕馀醒寐，遗浮沈沣浦。\par
有白鸥淡月，微波寄语，逍遥容与。\par

\section{六丑}
\textbf{彭元逊}\par\vspace{0.5em}
似东风老大，那复有、当时风气。\par
有情不收，江山身是寄。\par
浩荡何世。\par
但忆临官道，暂来不住，便出门千里。\par
痴心指望回风坠。\par
扇底相逢，钗头微缀。\par
他家万条千缕，解遮亭障驿，不隔江水。\par
瓜洲曾舣，等行人岁岁。\par
日下长秋，城乌夜起。\par
帐庐好在春睡。\par
共飞归湖上，草青无地。\par
雨、春心如腻。\par
欲待化、丰乐楼前，青门都废。\par
何人念、流落无几。\par
点点抟作，雪绵松润，为君泪。\par

\section{瑞鹤仙}
\textbf{陆淞}\par\vspace{0.5em}
脸霞红印枕。\par
睡觉来、冠儿还是不整。\par
屏间麝煤冷。\par
但眉峰压翠，泪珠弹粉。\par
堂深昼永。\par
燕交飞、风帘露井。\par
恨无人，与说相思，近日带围宽尽。\par
重省。\par
残灯朱幌，淡月纱窗，那时风景。\par
阳台路迥。\par
云雨梦，便无准。\par
待归来，先指花梢教看，却把心期细问。\par
问因循、过了青春，怎生意稳。\par

\section{减字浣溪沙・浣溪沙}
\textbf{贺铸}\par\vspace{0.5em}
楼角初销一缕霞。\par
淡黄杨柳暗栖鸦。\par
玉人和月摘梅花。\par
笑拈粉香归洞户，更垂帘幕护窗纱。\par
东风寒似夜来些。\par

\section{天门谣}
\textbf{贺铸}\par\vspace{0.5em}
牛渚天门险。\par
限南北、七雄豪占。\par
清雾敛。\par
与闲人登览。\par
待月上潮平波滟滟。\par
塞管轻吹新《阿滥》。\par
风满槛。\par
历历数、西州更点。\par

\section{石州引・石州慢}
\textbf{贺铸}\par\vspace{0.5em}
薄雨初寒，斜照弄晴，春意空阔。\par
长亭柳色缠黄，远客一枝先折。\par
烟横水际，映带几点归鸦，东风销尽龙沙雪。\par
还记初关来，恰而今时节。\par
将发。\par
画楼芳酒，红泪清歌，顿成清别。\par
已是经年，杳杳音尘多绝。\par
欲知方寸，共有几许清愁，芭蕉不展丁香结。\par
枉望断天涯，两厌厌风月。\par

\section{望湘人}
\textbf{贺铸}\par\vspace{0.5em}
厌莺声到枕，花气动帘，醉魂愁梦相半。\par
被惜馀薰，带惊剩眼。\par
几许伤春春晚。\par
泪竹痕鲜，佩兰香老，湘天浓暖。\par
记小江、风月佳时，屡约非烟游伴。\par
须信鸾弦易断。\par
奈云和再鼓，曲终人远。\par
认罗袜无踪，旧处弄波清浅。\par
青翰棹舣，白苹洲畔。\par
尽木临皋飞观。\par
不解寄、一字相思，幸有归来双燕。\par

\section{蝶恋花}
\textbf{贺铸}\par\vspace{0.5em}
几许伤春春复暮。\par
杨柳清阴，偏碍游丝度。\par
天际小山桃叶步。\par
白苹花满湔裙处。\par
竟日微吟长短句。\par
帘影灯昏，心寄胡琴语。\par
数点雨声风约住。\par
朦胧淡月云来去。\par

\section{夏云峰}
\textbf{仲殊}\par\vspace{0.5em}
天阔云高，溪横水远。\par
晚日寒生轻晕。\par
闲阶静、杨花渐少，朱门掩、莺声犹嫩。\par
悔匆匆、过却清明，旋占得馀芳，已成幽恨。\par
都几日阴沈，连宵慵困。\par
起来韶华都尽。\par
怨入双眉闲斗损。\par
乍品得情怀，看承全近。\par
深深态、无非自许。\par
厌厌意、终羞人间。\par
争知道、梦里蓬莱，待忘了馀香，时传音信。\par
纵留得莺花，东风不住，也则眼前愁闷。\par

\section{汉宫春}
\textbf{李邴}\par\vspace{0.5em}
潇洒江梅，向竹梢疏处，横两三枝。\par
东君也不爱惜，雪压霜欺。\par
无情燕子，怕春寒、轻失花期。\par
却是有，年年塞雁，归来曾见开时。\par
清浅小溪如练，问玉堂何似，茅舍疏篱。\par
伤心故人去後，冷落新诗。\par
微云淡月，对江天、分付他谁。\par
空自忆，清香未减，风流不在人知。\par

\section{贺新郎}
\textbf{潘汾}\par\vspace{0.5em}
篆缕销香鼎。\par
翠沈沈、庭阴转午，画堂人静。\par
芳草王孙知何处，惟有杨花糁径。\par
正玉枕、瞢腾初醒。\par
门外残红春已去，镇无聊、酒厌厌病。\par
云髻，未整。\par
江南旧事休重省。\par
但天涯、寻消问息，断鸿难倩。\par
月满西楼凭阑久，依旧归期未定。\par
便只恐、瓶沈金井。\par
嘶骑不来银烛暗，枉教人、立尽梧桐影。\par
谁伴我，对鸾镜。\par

\section{忆王孙}
\textbf{李重元}\par\vspace{0.5em}
萋萋芳草忆王孙。\par
柳外楼高空断魂。\par
杜宇声声不忍闻。\par
欲黄昏。\par
雨打梨花深闭门。\par

\section{贺新郎}
\textbf{李玉}\par\vspace{0.5em}
篆缕销金鼎。\par
醉沈沈、庭阴转午，画堂人静。\par
芳草王孙知何处，惟有杨花糁径。\par
渐玉枕、腾腾春醒。\par
帘外残红春已透，镇无聊、酒厌厌病。\par
云鬓乱，未整。\par
江南旧事休重省。\par
遍天涯、寻消问息，断鸿难倩。\par
月满西楼凭栏久，依旧归期未定。\par
又只恐、瓶沈金井。\par
嘶骑不来银烛暗，枉教人、立尽梧桐影。\par
谁伴我，对鸾镜。\par

\section{临江仙}
\textbf{陈与义}\par\vspace{0.5em}
高咏楚词酬午日，天涯节序匆匆。\par
榴花不似舞裙红。\par
无人知此意，歌罢满帘风。\par
万事一身伤老矣，戌葵凝笑墙东。\par
酒杯深浅去年同。\par
试浇桥下水，今夕到湘中。\par

\section{临江仙}
\textbf{陈与义}\par\vspace{0.5em}
忆昔午桥桥上饮，坐中多是豪英。\par
长沟流月去无声。\par
杏花疏影里，吹笛到天明。\par
二十馀年如一梦，此身虽在堪惊。\par
闲登小阁看新晴。\par
古今多少事，渔唱起三更。\par

\section{兰陵王}
\textbf{张元干}\par\vspace{0.5em}
卷朱箔。\par
朝雨轻阴乍阁。\par
阑干外，烟柳弄晴，芳草侵阶映红药。\par
东风妒花恶。\par
吹落。\par
梢头嫩萼。\par
屏山掩，沈水倦熏，中酒心情怕杯勺。\par
寻思旧京洛。\par
正年少疏狂，歌笑迷著。\par
障泥油壁催梳掠。\par
曾驰道同载，上林携手，灯夜初过早共约。\par
又争信漂泊。\par
寂寞。\par
念行乐。\par
甚粉淡衣襟，音断弦索。\par
琼枝璧月春如昨。\par
怅别後华表，那回双鹤。\par
相思除是，向醉里、暂忘却。\par

\section{石州慢}
\textbf{张元干}\par\vspace{0.5em}
寒水依痕，春意渐回，沙际烟阔。\par
溪梅晴照生香，冷芯数枝争发。\par
天涯旧恨，试看几许消魂，长亭门外山重叠。\par
不尽眼中青，是愁来时节。\par
情切。\par
画楼深闭，想见东风，暗销肌雪。\par
辜负枕前云雨，尊前花月。\par
心期切处，更有多少凄凉，殷勤留与归时说。\par
到得却相逢，恰经年离别。\par

\section{薄幸}
\textbf{吕渭老}\par\vspace{0.5em}
青楼春晚。\par
昼寂寂、梳匀又懒。\par
乍听得、鸦啼莺弄，惹起新愁无限。\par
记年时、偷掷春心，花间隔雾遥相见。\par
便角枕题诗，宝钗贳酒，共醉青苔深院。\par
怎忘得、回廊下，携手处、花明月满。\par
如今但暮雨，蜂愁蝶恨，小窗闲对芭蕉展。\par
却谁拘管。\par
尽无言、闲品秦筝，泪满参差雁。\par
腰支渐小，心与杨花共远。\par

\section{满江红}
\textbf{岳飞}\par\vspace{0.5em}
怒发冲冠，凭阑处、潇潇雨歇。\par
抬望眼、仰天长啸，壮怀激烈。\par
三十功名尘与土，八千里路云和月。\par
莫等闲、白了少年头，空悲切。\par
靖康耻，犹未雪。\par
臣子恨，何时灭。\par
驾长车踏破，贺兰山缺。\par
壮志饥餐胡虏肉，笑谈渴饮匈奴血。\par
待从头、收拾旧山河，朝天阙。\par

\section{水龙吟}
\textbf{陈亮}\par\vspace{0.5em}
闹花深处层楼，画帘半卷东风软。\par
春归翠陌，平莎茸嫩，垂杨金浅。\par
迟日催花，淡云阁雨，轻寒轻暖。\par
恨芳菲世界，游人未赏，都付与、莺和燕。\par
寂寞凭高念远。\par
向南楼、一声归雁。\par
金钗斗草，青丝勒马，风流云散。\par
罗绶分香，翠绡封泪，几多幽怨。\par
正销魂，又是疏烟淡月，子规声断。\par

\section{宴山亭・燕山亭}
\textbf{张}\par\vspace{0.5em}
幽梦初回，重阴未开，晓色吹成疏雨。\par
竹槛气寒，蕙畹声摇，新绿暗通南浦。\par
未有人行，才半启、回廊朱户。\par
无绪。\par
空望极霓旌，锦书难据。\par
苔径追忆曾游，念谁伴、秋千采绳芳柱。\par
犀奁黛卷，凤枕云孤，应也几番凝伫。\par
怎得伊来，花雾绕、小堂深处。\par
留住。\par
直到老、不教归去。\par

\section{唐多令}
\textbf{刘过}\par\vspace{0.5em}
芦叶满汀洲。\par
塞沙带浅流。\par
二十年、重过南楼。\par
柳下系舟犹未稳，能几日、又中秋。\par
黄鹤断矶头。\par
故人今不在。\par
旧江山，浑是新愁。\par
欲买桂花同载酒，终不似、少年游。\par

\section{点绛唇}
\textbf{姜夔}\par\vspace{0.5em}
燕雁无心，太湖西畔随云去。\par
数峰清苦。\par
商略黄昏雨。\par
第四桥边，拟共天随住。\par
今何许。\par
凭栏怀古。\par
残柳参差舞。\par

\section{鹧鸪天}
\textbf{姜夔}\par\vspace{0.5em}
肥水东流无尽期。\par
当初不合种相思。\par
梦中未比丹青见，暗里忽惊山鸟啼。\par
春未绿，鬓先丝。\par
人间别久不成悲。\par
谁教岁岁红莲夜，两处沈吟各自知。\par

\section{杏花天}
\textbf{姜夔}\par\vspace{0.5em}
绿丝低拂鸳鸯浦。\par
想桃叶、当时唤渡。\par
又将愁眼与春风，待去。\par
倚兰桡、更少驻。\par
金陵路。\par
莺吟燕舞。\par
算潮水、知人最苦。\par
满汀芳草不成归，日暮。\par
更移舟、向甚处。\par

\section{踏莎行}
\textbf{姜夔}\par\vspace{0.5em}
燕燕轻盈，莺莺娇软。\par
分明又向华胥见。\par
夜长争得薄情知，春初早被相思染。\par
别後书辞，别时针线。\par
离魂暗逐郎行远。\par
淮南皓月冷千山，冥冥归去无人管。\par

\section{霓裳中序第一}
\textbf{姜夔}\par\vspace{0.5em}
亭皋正望极。\par
乱落江莲归未得。\par
多病却无气力。\par
况纨扇渐疏，罗衣初萦。\par
流光过隙。\par
叹杏梁、双燕如客。\par
人何在，一帘淡月，仿佛照颜色。\par
幽寂。\par
乱蛩吟壁。\par
动庾信、清愁似织。\par
沈思年少浪迹。\par
笛里关山，柳下坊陌。\par
坠红无信息。\par
漫暗水，涓涓溜碧。\par
漂零久，而今何意，醉卧酒垆侧。\par

\section{庆宫春・高阳台}
\textbf{姜夔}\par\vspace{0.5em}
双桨莼波，一蓑松雨，暮愁渐满空阔。\par
呼我盟鸥，翩翩欲下，背人还过木末。\par
那回归去，荡云雪、孤舟夜发。\par
伤心重见，依约眉山，黛痕低压。\par
采香径里春寒，老子婆娑，自歌谁答。\par
垂虹西望，飘然引去，此兴平生难遏。\par
酒醒波远，政凝想、明素袜。\par
如今安在，唯有阑干，伴人一霎。\par

\section{齐天乐}
\textbf{姜夔}\par\vspace{0.5em}
庾郎先自吟愁赋。\par
凄凄更闻私语。\par
露湿铜铺，苔侵石井，都是曾听伊处。\par
哀音似诉。\par
正思妇无眠，起寻机杼。\par
曲曲屏山，夜凉独自甚情绪。\par
西窗又吹暗雨。\par
为谁频断续，相和砧杵。\par
候馆迎秋，离宫吊月，别有伤心无数。\par
豳诗漫与。\par
笑篱落呼灯，世间儿女。\par
写入琴丝，一声声更苦。\par

\section{一萼红}
\textbf{姜夔}\par\vspace{0.5em}
古城阴。\par
有官梅几许，红萼未宜簪。\par
池面冰胶，墙腰雪老，云意还又沈沈。\par
翠藤共、闲穿径竹，渐笑语、惊起卧沙禽。\par
野老林泉，故王台榭，呼唤登临。\par
南去北来何事，荡湘云楚水，目极伤心。\par
朱户黏鸡，金盘簇燕，空叹时序侵寻。\par
记曾共，西楼雅集，想垂杨、还袅万丝金。\par
待得归鞍到时，只怕春深。\par

\section{念奴娇}
\textbf{姜夔}\par\vspace{0.5em}
闹红一舸，记来时、尝与鸳鸯为侣。\par
三十六陂人未到，水佩风裳无数。\par
翠叶吹凉，玉容销酒，更洒菰蒲雨。\par
嫣然摇动，冷香飞上诗句。\par
日暮。\par
青盖亭亭，情人不见，争忍凌波去。\par
只恐舞衣寒易落，愁入西风南浦。\par
高柳垂阴，老鱼吹浪，留我花间住。\par
田田多少，几回沙际归路。\par

\section{琵琶仙・自度曲}
\textbf{姜夔}\par\vspace{0.5em}
双桨来时，有人似、旧曲桃根桃叶。\par
歌扇轻约飞花，蛾眉正奇绝。\par
春渐远、汀洲自绿，更添了、几声啼。\par
十里扬州，三生杜牧，前事休说。\par
又还是、宫烛分烟，奈愁里、匆匆换时节。\par
都把一襟芳思，与空阶榆荚。\par
千万缕、藏鸦细柳，为玉尊、起舞回雪。\par
想见西出阳关，故人初别。\par

\section{八归}
\textbf{姜夔}\par\vspace{0.5em}
芳莲坠粉，蔬桐吹绿，庭院暗雨乍歇。\par
无端抱影销魂处，还见筱墙萤暗，藓阶蛩切。\par
送客重寻西去路，问水面、琵琶谁拨。\par
最可惜、一片江山，总付与啼。\par
长恨相从未款，而今何事，又对西风离别。\par
渚寒烟淡，棹移人远，缥缈行舟如叶。\par
想文君望久，倚竹愁生步罗袜。\par
归来後、翠尊双饮，下了珠帘，玲珑闲看月。\par

\section{扬州慢}
\textbf{姜夔}\par\vspace{0.5em}
淮左名都，竹西佳处，解鞍少驻初程。\par
过春风十里，尽荠麦青青。\par
自胡马窥江去後，废池乔木，犹厌言兵。\par
渐黄昏，清角吹寒，都在空城。\par
杜郎俊赏，算而今，重到须惊。\par
纵豆蔻词工，青楼梦好，难赋深情。\par
二十四桥仍在，波心荡、冷月无声。\par
念桥边红药，年年知为谁生。\par

\section{长亭怨慢}
\textbf{姜夔}\par\vspace{0.5em}
渐吹尽、枝头香絮。\par
是处人家，绿深门户。\par
远浦萦回，暮帆零乱向何许。\par
阅人多矣，谁得似、长亭树。\par
树若有情时，不会得、青青如此。\par
日暮。\par
望高城不见，只见乱山无数。\par
韦郎去也，怎忘得、玉环分付。\par
第一是、早早归来，怕红萼、无人为主。\par
算空有并刀，难翦离愁千缕。\par

\section{淡黄柳}
\textbf{姜夔}\par\vspace{0.5em}
空城晓月。\par
吹入垂杨陌。\par
马上单衣寒恻恻。\par
看尽鹅黄嫩绿，都是江南旧相识。\par
正岑寂。\par
明朝又寒食。\par
强携酒、小桥宅，怕梨花落尽成秋色。\par
燕燕飞来，问春何在，唯有池塘自碧。\par

\section{暗香}
\textbf{姜夔}\par\vspace{0.5em}
旧时月色。\par
算几番照我，梅边吹笛。\par
唤起玉人，不管清寒与攀摘。\par
何逊而今渐老，都忘却、春风词笔。\par
但怪得、竹外疏花，香冷入瑶席。\par
江国。\par
正寂寂。\par
叹寄与路遥，夜雪初积。\par
翠尊易泣。\par
红萼无言耿相忆。\par
长记曾携手处，千树压、西湖寒碧。\par
又片片、吹尽也，几时见得。\par

\section{疏影}
\textbf{姜夔}\par\vspace{0.5em}
苔枝缀玉。\par
有翠禽小小，枝上同宿。\par
客里相逢，篱角黄昏，无言自倚修竹。\par
昭君不惯胡沙远，但暗忆、江南江北。\par
想佩环、月夜归来，化作此花幽独。\par
犹记深宫旧事，那人正睡里，飞近蛾绿。\par
莫似春风，不管盈盈，早与安排金屋。\par
还教一片随波去，又却怨、玉龙哀曲。\par
等恁时、重觅幽香，已入小窗横幅。\par

\section{翠楼吟}
\textbf{姜夔}\par\vspace{0.5em}
月冷龙沙，尘轻虎落，今年汉初赐。\par
新翻胡部曲，听毡幕、元戎歌吹。\par
层楼高峙。\par
看栏曲萦红，檐牙飞翠。\par
人姝丽。\par
粉香吹下，夜寒风细。\par
此地。\par
宜有词仙，拥素云黄鹤，与君游戏。\par
玉梯凝望久，叹芳草、萋萋千里。\par
天涯情味。\par
仗酒清愁，花销英气。\par
西山外。\par
晚来还卷，一帘秋霁。\par

\section{木兰花・玉楼春}
\textbf{钱惟演}\par\vspace{0.5em}
城上风光莺语乱。\par
城下烟波春拍岸。\par
绿杨芳草几时休，泪眼愁肠先已断。\par
情怀渐变成衰晚。\par
鸾鉴朱颜惊暗换。\par
昔年多病厌芳尊，今日芳尊惟恐浅。\par

\section{苏幕遮}
\textbf{范仲淹}\par\vspace{0.5em}
碧云天，黄叶地。\par
秋色连波，波上寒烟翠。\par
山映斜阳天接水。\par
芳草无情，更在斜阳外。\par
黯乡魂，追旅思。\par
夜夜除非，好梦留人睡。\par
明月楼高休独倚。\par
酒入愁肠，化作相思泪。\par

\section{御街行}
\textbf{范仲淹}\par\vspace{0.5em}
纷纷堕叶飘香砌。\par
夜寂静、寒声碎。\par
真珠帘卷玉楼空，天淡银河垂地。\par
年年今夜，月华如练，长是人千里。\par
愁肠已断无由醉。\par
酒未到、先成泪。\par
残灯明灭枕头欹。\par
谙尽孤眠滋味。\par
都来此事，眉间心上，无计相回避。\par

\section{曲玉管}
\textbf{柳永}\par\vspace{0.5em}
陇首云飞，江边日晚，烟波满目凭阑久。\par
立望关河萧索，千里清秋。\par
忍凝眸。\par
杳杳神京，盈盈仙子，别来锦字终难偶。\par
断雁无凭，冉冉飞下汀洲。\par
思悠悠。\par
暗想当初，有多少、幽欢佳会，岂知聚散难期，翻成雨恨云愁。\par
阻追游。\par
每登山临水，惹起平生心事，一场消黯，永日无言，却下层楼。\par

\section{雨霖铃}
\textbf{柳永}\par\vspace{0.5em}
寒蝉凄切。\par
对长亭晚，骤雨初歇。\par
都门帐饮无绪，留恋处、兰舟催发。\par
执手相看泪眼，竟无语凝噎。\par
念去去、千里烟波，暮霭沈沈楚天阔。\par
多情自古伤离别。\par
更那堪、冷落清秋节。\par
今宵酒醒何处，杨柳岸、晓风残月。\par
此去经年，应是良辰、好景虚设。\par
便纵有、千种风情，更与何人说。\par

\section{采莲令}
\textbf{柳永}\par\vspace{0.5em}
月华收，云淡霜天曙。\par
西征客、此时情苦。\par
翠娥执手送临歧，轧轧开朱户。\par
千娇面、盈盈伫立，无言有泪，断肠争忍回顾。\par
一叶兰舟，便恁急桨凌波去。\par
贪行色、岂知离绪。\par
万般方寸，但饮恨，脉脉同谁语。\par
更回首、重城不见，寒江天外，隐隐两三烟树。\par

\section{凤栖梧・蝶恋花}
\textbf{柳永}\par\vspace{0.5em}
伫倚危楼风细细。\par
望极春愁，黯黯生天际。\par
草色烟光残照里。\par
无言谁会凭阑意。\par
拟把疏狂图一醉。\par
对酒当歌，强乐还无味。\par
衣带渐宽终不悔。\par
为伊消得人憔悴。\par

\section{浪淘沙}
\textbf{柳永}\par\vspace{0.5em}
梦觉、透窗风一线，寒灯吹息。\par
那堪酒醒，又闻空阶，夜雨频滴。\par
嗟因循、久作天涯客。\par
负佳人、几许盟言，便忍把、从前欢会，陡顿翻成忧戚。\par
愁极。\par
再三追思，洞房深处，几度饮散歌阑，香暖鸳鸯被，岂暂时疏散，费伊心力。\par
云尤雨，有万般千种，相怜相惜。\par
恰到如今，天长漏永，无端自家疏隔。\par
知何时、却拥秦云态，愿低帏昵枕，轻轻细说与，江乡夜夜，数寒更思忆。\par

\section{定风波}
\textbf{柳永}\par\vspace{0.5em}
自春来、惨绿愁红，芳心是事可可。\par
日上花梢，莺穿柳带，犹压香衾卧。\par
暖酥消，腻云。\par
终日厌厌倦梳裹。\par
无那。\par
恨薄情一去，音书无个。\par
早知恁麽。\par
悔当初、不把雕鞍锁。\par
向鸡窗、只与蛮笺象管，拘束教吟课。\par
镇相随，莫抛躲。\par
针线闲拈伴伊坐。\par
和我。\par
免使年少，光阴虚过。\par

\section{少年游}
\textbf{柳永}\par\vspace{0.5em}
长安古道马迟迟。\par
高柳乱蝉栖。\par
夕阳岛外，秋风原上，目断四天垂。\par
归云一去无踪迹，何处是前期。\par
狎兴生疏，酒徒萧索，不似去年时。\par

\section{戚氏}
\textbf{柳永}\par\vspace{0.5em}
晚秋天。\par
一霎微雨洒庭轩。\par
槛菊萧疏，井梧零乱惹残烟。\par
凄然。\par
望江关。\par
飞云黯淡夕阳间。\par
当时宋玉悲感，向此临水与登山。\par
远道迢递，行人凄楚，倦听陇水潺。\par
正蝉吟败叶，蛩响衰草，相应喧喧。\par
孤馆度日如年。\par
风露渐变，悄悄至更阑。\par
长天净，绛河清浅，皓月婵娟。\par
思绵绵。\par
夜永对景，那堪屈指，暗想从前。\par
未名未禄，绮陌红楼，往往经岁迁延。\par
帝里风光好，当年少日，暮宴朝欢。\par
况有狂朋怪侣，遇当歌、对酒竞留连。\par
别来迅景如梭，旧游似梦，烟水程何限。\par
念利名、憔悴长萦绊。\par
追往事、空惨愁颜。\par
漏箭移、稍觉轻寒。\par
渐呜咽、画角数声残。\par
对闲窗畔，停灯向晓，抱影无眠。\par

\section{夜半乐}
\textbf{柳永}\par\vspace{0.5em}
冻云黯淡天气，扁舟一叶，乘兴离江渚。\par
渡万壑千岩，越溪深处。\par
怒涛渐息，樵风乍起，更闻商旅相呼。\par
片帆高举。\par
泛画鹢、翩翩过南浦。\par
望中酒旆闪闪，一簇烟村，数行霜树。\par
残日下，渔人鸣榔归去。\par
败荷零落，衰杨掩映，岸边两两三三，浣沙游女。\par
避行客、含羞笑相语。\par
到此因念，绣阁轻抛，浪萍难驻。\par
叹後约丁宁竟何据。\par
惨离怀，空恨岁晚归期阻。\par
凝泪眼、杳杳神京路。\par
断鸿声远长天暮。\par

\section{玉蝴蝶}
\textbf{柳永}\par\vspace{0.5em}
望处雨收云断，凭阑悄悄，目送秋光。\par
晚景萧疏，堪动宋玉悲凉。\par
水风轻、苹花渐老，月露冷、梧叶飘黄。\par
遣情伤。\par
故人何在，烟水茫茫。\par
难忘。\par
文期酒会，几孤风月，屡变星霜。\par
海阔山遥，未知何处是潇湘。\par
念双燕、难凭远信，指暮天、空识归航。\par
黯相望。\par
断鸿声里，立尽斜阳。\par

\section{八声甘州}
\textbf{柳永}\par\vspace{0.5em}
对潇潇、暮雨洒江天，一番洗清秋。\par
渐霜风凄惨，关河冷落，残照当楼。\par
是处红衰翠减，苒苒物华休。\par
惟有长江水，无语东流。\par
不忍登高临远，望故乡渺邈，归思难收。\par
叹年来踪迹，何事苦淹留。\par
想佳人、妆楼望，误几回、天际识归舟。\par
争知我、倚阑干处，正恁凝愁。\par

\section{竹马子・竹马儿}
\textbf{柳永}\par\vspace{0.5em}
登孤垒荒凉，危亭旷望，静临烟渚。\par
对雌霓挂雨，雄风拂槛，微收烦暑。\par
渐觉一叶惊秋，残蝉噪晚，素商时序。\par
览景想前欢，指神京，非雾非烟深处。\par
向此成追感，新愁易积，故人难聚。\par
凭高尽日凝伫。\par
赢得消魂无语。\par
极目霁霭霏微，暝鸦零乱，萧索江城暮。\par
南楼画角，又送残阳去。\par

\section{迷神引}
\textbf{柳永}\par\vspace{0.5em}
一叶扁舟轻帆卷。\par
暂泊楚江南岸。\par
孤城暮角，引胡茄怨。\par
水茫茫，平沙雁、旋惊散。\par
烟敛寒林簇，画屏展。\par
天际遥山小，黛眉浅。\par
旧赏轻抛，到此成游宦。\par
觉客程劳，年光晚。\par
异乡风物，忍萧索、当愁眼。\par
帝城赊，秦楼阻，旅魂乱。\par
芳草连空阔，残照满。\par
佳人无消息，断云远。\par

\section{醉垂鞭}
\textbf{张先}\par\vspace{0.5em}
双蝶绣罗裙。\par
东池宴。\par
初相见。\par
朱粉不深匀。\par
闲花淡淡春。\par
细看诸处好。\par
人人道。\par
柳腰身。\par
昨日乱山昏。\par
来时衣上云。\par

\section{一丛花令・一丛花}
\textbf{张先}\par\vspace{0.5em}
伤高怀远几时穷。\par
无物似情浓。\par
离愁正引千丝乱，更东陌、飞絮。\par
嘶骑渐遥，征尘不断，何处认郎踪。\par
双鸳池沼水溶溶。\par
南北小桡通。\par
梯横画阁黄昏後，又还是、斜月帘栊。\par
沈恨细思，不如桃杏，犹解嫁东风。\par

\section{天仙子}
\textbf{张先}\par\vspace{0.5em}
水调数声持酒听。\par
午醉醒来愁未醒。\par
送春春去几时回，临晚镜。\par
伤流景。\par
往事後期空记省。\par
沙上并禽池上暝。\par
云破月来花弄影。\par
重重帘幕密遮灯，风不定。\par
人初静。\par
明日落红应满径。\par

\section{千秋岁}
\textbf{张先}\par\vspace{0.5em}
数声鶗鴂。\par
又报芳菲歇。\par
惜春更把残红折。\par
雨轻风色暴，梅子青时节。\par
永丰柳，无人尽日飞花雪。\par
莫把幺弦拨。\par
怨极弦能说。\par
天不老，情难绝。\par
心似双丝网，中有千千结。\par
夜过也，东窗未白凝残月。\par

\section{青门引・青门饮}
\textbf{张先}\par\vspace{0.5em}
乍暖还轻冷。\par
风雨晚来方定。\par
庭轩寂寞近清明，残花中酒，又是去年病。\par
楼头画角风吹醒。\par
入夜重门静。\par
那堪更被明月，隔墙送过秋千影。\par

\section{浣溪沙}
\textbf{晏殊}\par\vspace{0.5em}
一曲新词酒一杯。\par
去年天气旧亭台。\par
夕阳西下几时回。\par
无可奈何花落去，似曾相识燕归来。\par
小园香径独徘徊。\par

\section{浣溪沙}
\textbf{晏殊}\par\vspace{0.5em}
一向年光有限身。\par
等闲离别易消魂。\par
酒筵歌席莫辞频。\par
满目山河空念远，落花风雨更伤春。\par
不如怜取眼前人。\par

\section{清平乐}
\textbf{晏殊}\par\vspace{0.5em}
金风细细。\par
叶叶梧桐坠。\par
绿酒初尝人易醉。\par
一枕小窗浓睡。\par
紫微朱槿花残。\par
斜阳却照阑干。\par
双燕欲归时节，银屏昨夜微寒。\par

\section{清平乐}
\textbf{晏殊}\par\vspace{0.5em}
红笺小字。\par
说尽平生意。\par
鸿雁在云鱼在水。\par
惆怅此情难寄。\par
斜阳独倚西楼。\par
遥山恰对帘钩。\par
人面不知何处，绿波依旧东流。\par

\section{木兰花・玉楼春}
\textbf{晏殊}\par\vspace{0.5em}
燕鸿过後莺归去。\par
细算浮生千万绪。\par
长於春梦几多时，散似秋云无觅处。\par
闻琴解佩神仙侣。\par
挽断罗衣留不住。\par
劝君莫作独醒人，烂醉花间应有数。\par

\section{木兰花・玉楼春}
\textbf{晏殊}\par\vspace{0.5em}
池塘水绿风微暖。\par
记得玉真初见面。\par
重头歌韵响铮琮，入破舞腰红乱旋。\par
玉钩阑下香阶畔。\par
醉後不知斜日晚。\par
当时共我赏花人，点检如今无一半。\par

\section{踏莎行}
\textbf{晏殊}\par\vspace{0.5em}
祖席离歌，长亭别宴。\par
香尘已隔犹回面。\par
居人匹马映林嘶，行人去棹依波转。\par
画阁魂消，高楼目断。\par
斜阳只送平波远。\par
无穷无尽是离愁，天涯地角寻思遍。\par

\section{踏莎行}
\textbf{晏殊}\par\vspace{0.5em}
小径红稀，芳郊绿遍。\par
高台树色阴阴见。\par
春风不解禁杨花，乱扑行人面。\par
翠叶藏莺，朱帘隔燕。\par
炉香静逐游丝转。\par
一场愁梦酒醒时，斜阳却照深深院。\par

\section{玉楼春}
\textbf{晏殊}\par\vspace{0.5em}
绿杨芳草长亭路。\par
年少抛人容易去。\par
楼头残梦五更钟，花底离情三月雨。\par
无情不似多情苦。\par
一寸还成千万缕。\par
天涯地角有穷时，只有相思无尽处。\par

\section{玉楼春}
\textbf{宋祁}\par\vspace{0.5em}
东城渐觉风光好。\par
皱波纹迎客棹。\par
绿杨烟外晓寒轻，红杏枝头春意闹。\par
浮生长恨欢娱少。\par
肯爱千金轻一笑。\par
为君持酒劝斜阳，且向花间留晚照。\par

\section{采桑子}
\textbf{欧阳修}\par\vspace{0.5em}
群芳过後西湖好，狼籍残红。\par
飞絮蒙蒙。\par
垂柳阑干尽日风。\par
笙歌散尽游人去，始觉春空。\par
垂下帘栊。\par
双燕归来细雨中。\par

\section{诉衷情}
\textbf{欧阳修}\par\vspace{0.5em}
清晨帘幕卷轻霜。\par
呵手试梅妆。\par
都缘自有离恨，故画作远山长。\par
思往事，惜流芳。\par
易成伤。\par
拟歌先敛，欲笑还颦，最断人肠。\par

\section{踏莎行}
\textbf{欧阳修}\par\vspace{0.5em}
候馆梅残，溪桥柳细。\par
草熏风暖摇征辔。\par
离愁渐远渐无穷，迢迢不断如春水。\par
寸寸柔肠，盈盈粉泪。\par
楼高莫近危阑倚。\par
平芜尽处是春山，行人更在春山外。\par

\section{蝶恋花}
\textbf{欧阳修}\par\vspace{0.5em}
独倚危楼风细细。\par
望极离愁，黯黯生天际。\par
草色山光残照里。\par
无人会得凭阑意。\par
也拟疏狂图一醉。\par
对酒当歌，强饮还无味。\par
衣带渐宽都不悔。\par
况伊销得人憔悴。\par

\section{玉楼春}
\textbf{欧阳修}\par\vspace{0.5em}
池塘水绿春微暖。\par
记得玉真初见面。\par
从头歌韵响铮鏦，入破舞腰红乱旋。\par
玉钩帘下香阶畔。\par
醉後不知红日晚。\par
当时共我赏花人，点检如今无一半。\par

\section{玉楼春}
\textbf{欧阳修}\par\vspace{0.5em}
燕鸿过後春归去。\par
细算浮生千万绪。\par
来如春梦几多时，去似朝云无觅处。\par
闻琴解神仙侣。\par
挽断罗衣留不住。\par
劝君莫作独醒人，烂醉花间应有数。\par

\section{玉楼春}
\textbf{欧阳修}\par\vspace{0.5em}
别後不知君远近。\par
触目凄凉多少闷。\par
渐行渐远渐无书，水阔鱼沈何处问。\par
夜深风竹敲秋韵。\par
万叶千声皆是恨。\par
故欹单枕梦中寻，梦又不成灯又尽。\par

\section{浪淘沙}
\textbf{欧阳修}\par\vspace{0.5em}
把酒祝东风。\par
且共从容。\par
垂杨紫陌洛城东。\par
总是当时携手处，游遍芳丛。\par
聚散苦匆匆。\par
此恨无穷。\par
今年花胜去年红。\par
可惜明年花更好，知与谁同。\par

\section{风入松}
\textbf{俞国宝}\par\vspace{0.5em}
一春长费买花钱。\par
日日醉花边。\par
玉骢惯识西湖路，骄嘶过、沽酒垆前。\par
红杏香中箫鼓，绿杨影里秋千。\par
暖风十里丽人天。\par
花压髻云偏。\par
画船载取春归去，馀情寄、湖水湖烟。\par
明日重扶残醉，来寻陌上花钿。\par

\section{绮罗香}
\textbf{史达祖}\par\vspace{0.5em}
做冷欺花，将烟困柳，千里偷催春暮。\par
尽日冥迷，愁里欲飞还住。\par
惊粉重、蝶宿西园，喜泥润、燕归南浦。\par
最妨它、佳约风流，钿车不到杜陵路。\par
沈沈江上望极，还被春潮晚急，难寻官渡。\par
隐约遥峰，和泪谢娘眉妩。\par
临断岸、新绿生时，是落红、带愁流处。\par
记当日、门掩梨花，翦灯深夜语。\par

\section{双双燕}
\textbf{史达祖}\par\vspace{0.5em}
过春社了，度帘幕中间，去年尘冷。\par
差池欲住，试入旧巢相并。\par
还相雕梁藻井。\par
又软语、商量不定。\par
飘然快拂花梢，翠尾分开红影。\par
芳径。\par
芹泥雨润。\par
爱贴地争飞，竞夸轻俊。\par
红楼归晚，看足柳昏花暝。\par
应自栖香正稳。\par
便忘了、天涯芳信。\par
愁损翠黛双蛾，日日画阑独凭。\par

\section{东风第一枝}
\textbf{史达祖}\par\vspace{0.5em}
巧沁兰心，偷黏草甲，东风欲障新暖。\par
谩凝碧瓦难留，信知暮寒轻浅。\par
行天入镜，做弄出、轻松纤软。\par
料故园、不卷重帘，误了乍来双燕。\par
青未了、柳回白眼。\par
红欲断、杏开素面。\par
旧游忆著山阴，厚盟遂妨上苑。\par
寒炉重暖，便放慢春衫针线。\par
恐凤靴、挑菜归来，万一灞桥相见。\par

\section{喜迁莺}
\textbf{史达祖}\par\vspace{0.5em}
月波疑滴。\par
望玉壶天近，了无尘隔。\par
翠眼圈花，冰丝织练，黄道宝光相直。\par
自怜诗酒瘦，难应接、许多春色。\par
最无赖，是随香趁烛，曾伴狂客。\par
踪迹。\par
谩记忆。\par
老了杜郎，忍听东风笛。\par
柳院灯疏，梅厅雪在，谁与细倾春碧。\par
旧情拘未定，犹自学、当年游历。\par
怕万一，误玉人、夜寒帘隙。\par

\section{三姝媚}
\textbf{史达祖}\par\vspace{0.5em}
烟光摇缥瓦。\par
望晴檐多风，柳花如洒。\par
锦瑟横床，想泪痕尘影，凤弦常下。\par
倦出犀帷，频梦见、王孙骄马。\par
讳道相思，偷理绡裙，自惊腰衩。\par
惆怅南楼遥夜。\par
记翠箔张灯，枕肩歌罢。\par
又入铜驼，遍旧家门巷，首询声价。\par
可惜东风，将恨与、闲花俱谢。\par
记取崔徽模样，归来暗写。\par

\section{夜合花}
\textbf{史达祖}\par\vspace{0.5em}
柳锁莺魂，花翻蝶梦，自知愁染潘郎。\par
轻衫未揽，犹将泪点偷藏。\par
念前事，怯流光。\par
早春窥、酥雨池塘。\par
向销凝里，梅开半面，情满徐妆。\par
风丝一寸柔肠。\par
曾在歌边惹恨，烛底萦香。\par
芳机瑞锦，如何未织鸳鸯。\par
人扶醉，月依墙。\par
是当初、谁敢疏狂。\par
把闲言语，花房夜久，各自思量。\par

\section{八归}
\textbf{史达祖}\par\vspace{0.5em}
秋江带雨，寒沙萦水，人瞰画阁愁独。\par
烟蓑散响惊诗思，还被乱鸥飞去，秀句难续。\par
冷眼尽归图画上，认隔岸、微茫云屋。\par
想半属、渔市樵村，欲暮竞然竹。\par
须信风流未老，凭持酒、慰此凄凉心目。\par
一鞭南陌，几篙官渡，赖有歌眉舒绿。\par
只匆匆眺远，早觉闲愁挂乔木。\par
应难奈，故人天际，望彻淮山，相思无雁足。\par

\section{玉胡蝶・玉蝴蝶}
\textbf{史达祖}\par\vspace{0.5em}
晚雨未摧宫树，可怜闲叶，犹抱凉蝉。\par
短景归秋，吟思又接愁边。\par
漏初长、梦魂难禁，人渐老、风月俱寒。\par
想幽欢。\par
土花庭，虫网阑干。\par
无端。\par
啼蛄搅夜，恨随团扇，苦近秋莲。\par
一笛当楼，谢娘悬泪立风前。\par
故园晚、强留诗酒，新雁远、不致寒暄。\par
隔苍烟。\par
楚香罗袖，谁伴婵娟。\par

\section{秋霁}
\textbf{史达祖}\par\vspace{0.5em}
江水苍苍，望倦柳愁荷，共感秋色。\par
废阁先凉，古帘空暮，雁程最嫌风力。\par
故园信息。\par
爱渠入眼南山碧。\par
念上国。\par
谁是、脍鲈江汉未归客。\par
还又岁晚，瘦骨临风，夜闻秋声，吹动岑寂。\par
露蛩悲、清灯冷屋，翻书愁上鬓毛白。\par
年少俊游浑断得。\par
但可怜处，无奈苒苒魂惊，采香南浦，翦梅烟驿。\par

\section{宴清都}
\textbf{卢祖皋}\par\vspace{0.5em}
春讯飞琼管。\par
风日薄、度墙啼鸟声乱。\par
江城次第，笙歌翠合，绮罗香暖。\par
溶溶涧渌冰泮。\par
醉梦里、年华暗换。\par
料黛眉重锁隋堤，芳心还动梁苑。\par
新来雁阔云音，鸾分槛影，无计重见。\par
啼春细雨，笼愁澹月，恁时庭院。\par
离肠未语先断。\par
算犹有、凭高望眼。\par
更那堪、芳草连天，飞梅弄晚。\par

\section{江城子}
\textbf{卢祖皋}\par\vspace{0.5em}
画楼帘幕卷新晴。\par
掩银屏。\par
晓寒轻。\par
坠粉飘香，日日唤愁生。\par
暗数十年湖上路，能几度，著娉婷。\par
年华空自感飘零。\par
拥春酲。\par
对谁醒。\par
天阔云闲，无处觅箫声。\par
载酒买花年少事，浑不似，旧心情。\par

\section{小重山}
\textbf{章良能}\par\vspace{0.5em}
柳暗花明春事深。\par
小阑红芍药，已抽簪。\par
雨馀风软碎鸣禽。\par
迟迟日，犹带一分阴。\par
往事莫沈吟。\par
身闲时序好，且登临。\par
旧游无处不堪寻。\par
无寻处，惟有少年心。\par

\section{满庭芳}
\textbf{张}\par\vspace{0.5em}
月洗高梧，露幽草，宝钗楼外秋深。\par
土花沿翠，萤火坠墙阴。\par
静听寒声断续，微韵转、凄咽悲沈。\par
争求侣，殷勤劝织，促破晓机心。\par
儿时，曾记得，呼灯灌穴，敛步随音。\par
任满身花影，犹自追寻。\par
携向花堂戏斗，亭台小、笼巧妆金。\par
今休说，从渠床下，凉夜伴孤吟。\par

\section{花犯}
\textbf{吴文英}\par\vspace{0.5em}
小娉婷，清铅素靥，蜂黄暗偷晕。\par
翠翘敧鬓。\par
昨夜冷中庭，月下相认。\par
睡浓更苦凄风紧。\par
惊回心未稳。\par
送晓色、一壶葱茜，才知花梦准。\par
湘娥化作此幽芳，凌波路，古岸云沙遗恨。\par
临砌影，寒香乱、冻梅藏韵。\par
熏炉畔、旋移傍枕，还又见、玉人垂绀鬒。\par
料唤赏、清华池馆，台杯须满引。\par

\section{浣溪沙}
\textbf{吴文英}\par\vspace{0.5em}
门隔花深梦旧游。\par
夕阳无语燕归愁。\par
玉纤香动小帘钩。\par
落絮无声春堕泪，行云有影月含羞。\par
东风临夜冷於秋。\par

\section{浣溪沙}
\textbf{吴文英}\par\vspace{0.5em}
波面铜花冷不收。\par
玉人垂钓理纤钩。\par
月明池阁夜来秋。\par
江燕话归成晓别，水花红减似春休。\par
西风梧井叶先愁。\par

\section{点绛唇}
\textbf{吴文英}\par\vspace{0.5em}
卷尽愁云，素娥临夜新梳洗。\par
暗尘不起。\par
酥润凌波地。\par
辇路重来，仿佛灯前事。\par
情如水。\par
小楼熏被。\par
春梦笙歌里。\par

\section{祝英台近}
\textbf{吴文英}\par\vspace{0.5em}
采幽香，巡古苑，竹冷翠微路。\par
斗草溪根，沙印小莲步。\par
自怜两鬓清霜，一年寒食，又身在、云山深处。\par
昼闲度。\par
因甚天也悭春，轻阴便成雨。\par
绿暗长亭，归梦趁风絮。\par
有情花影阑干，莺声门径，解留我、霎时凝伫。\par

\section{祝英台近}
\textbf{吴文英}\par\vspace{0.5em}
翦红情，裁绿意，花信上钗股。\par
残日东风，不放岁华去。\par
有人添烛西窗，不眠侵晓，笑声转、新年莺语。\par
旧尊俎。\par
玉纤曾擘黄柑，柔香系幽素。\par
归梦湖边，还迷镜中路。\par
可怜千点吴霜，寒销不尽，又相对、落梅如雨。\par

\section{澡兰香}
\textbf{吴文英}\par\vspace{0.5em}
盘丝系腕，巧篆垂簪，玉隐绀纱睡觉。\par
银瓶露井，彩云窗，往事少年依约。\par
为当时、曾写榴裙，伤心红绡褪萼。\par
黍梦光阴渐老，汀洲烟箬。\par
莫唱江南古调，怨抑难招，楚江沈魄。\par
薰风燕乳，暗雨梅黄，午镜澡兰帘幕。\par
念秦楼、也拟人归，应翦菖浦自酌。\par
但怅望、一屡新蟾，随人天角。\par

\section{燕山亭}
\textbf{赵佶}\par\vspace{0.5em}
裁翦冰绡，打叠数重，冷淡燕脂匀注。\par
新样靓妆，艳溢香融，羞杀蕊珠宫女。\par
易得凋零，更多少、无情风雨。\par
愁苦。\par
闲院落凄凉，几番春暮。\par
凭寄离恨重重，这双燕，何曾会人言语。\par
天遥地远，万水千山，知他故宫何处。\par
怎不思量，除梦里、有时会去。\par
无据。\par
和梦也、有时不做。\par

\section{烛影摇红}
\textbf{廖世美}\par\vspace{0.5em}
霭霭春空，画楼森耸凌云渚。\par
紫薇登览最关情，绝妙夸能赋。\par
惆怅相思迟暮。\par
记当日、朱阑共语。\par
塞鸿难问，岸柳何穷，别愁纷絮。\par
催促年光，旧来流水知何处。\par
断肠何必更残阳，极目伤平楚。\par
晚霁波声带雨。\par
悄无人、舟横野渡。\par
数峰江上，芳草天涯，参差烟树。\par

\section{凤凰台上忆吹箫}
\textbf{李清照}\par\vspace{0.5em}
香冷金猊，被翻红浪，起来人未梳头。\par
任宝奁闲掩，日上帘钩。\par
生怕闲愁暗恨，多少事、欲说还休。\par
今年瘦，非干病酒，不是悲秋。\par
明朝，这回去也，千万遍阳关，也即难留。\par
念武陵春晚，云锁重楼。\par
记取楼前绿水，应念我、终日凝眸。\par
凝眸处，从今更数，几段新愁。\par

\section{醉花阴}
\textbf{李清照}\par\vspace{0.5em}
薄雾浓云愁永昼。\par
瑞脑消金兽。\par
佳节又重阳，玉枕纱厨，半夜凉初透。\par
东篱把洒黄昏後，有暗香盈袖。\par
莫道不消魂，帘卷西风，人似黄花瘦。\par

\section{念奴娇}
\textbf{李清照}\par\vspace{0.5em}
萧条庭院，有斜风细雨，重门须闭。\par
宠柳娇花寒食近，种种恼人天气。\par
险韵诗成，扶头酒醒，别是闲滋味。\par
征鸿过尽，万千心事难寄。\par
楼上几日春寒，帘垂四面，玉阑干慵倚。\par
被冷香消新梦觉，不许愁人不起。\par
清露晨流，新桐初引，多少游春意。\par
日高烟敛，更看今日晴未。\par

\section{永遇乐}
\textbf{李清照}\par\vspace{0.5em}
落日熔金，暮云合璧，人在何处。\par
染柳烟浓，吹梅笛怨，春意知几许。\par
元宵佳节，融和天气，次第岂无风雨。\par
来相召、香车宝马，谢他酒朋诗侣。\par
中州盛日，闺门多暇，记得偏重三五。\par
铺翠冠儿，拈金雪柳，簇带争济楚。\par
如今憔悴，风鬟霜鬓，怕见夜间出去。\par
不如向、帘儿底下，听人笑语。\par

\section{声声慢}
\textbf{李清照}\par\vspace{0.5em}
寻寻觅觅，冷冷清清，凄凄惨惨戚戚。\par
乍暖还寒时候，最难将息。\par
三杯两盏淡酒，怎敌他、晚来风急。\par
雁过也，正伤心，却是旧时相识。\par
满地黄花堆积。\par
憔悴损，如今有谁摘。\par
守著窗儿，独自怎生得黑。\par
梧桐更兼细雨，到黄昏、点点滴滴。\par
这次第，怎一个、愁字了得。\par

\section{苏武慢}
\textbf{蔡伸}\par\vspace{0.5em}
雁落平沙，烟笼寒水，古垒鸣笳声断。\par
青山隐隐，败叶萧萧，天际暝鸦零乱。\par
楼上黄昏，片帆千里归程，年华将晚。\par
望碧云空暮，佳人何处，梦魂俱远。\par
忆旧游、邃馆朱扉，小园香径，尚想桃花人面。\par
书盈锦轴，恨满金徽，难写寸心幽怨。\par
两地离愁，一尊芳酒，凄凉危阑倚遍。\par
尽迟留、凭仗西风，吹乾泪眼。\par

\section{柳梢青}
\textbf{蔡伸}\par\vspace{0.5em}
数声鶗鴂。\par
可怜又是，春归时节。\par
满院东风，海棠铺绣，梨花飘雪。\par
丁香露泣残枝，算未比、愁肠寸结。\par
自是休文，多情多感，不干风月。\par

\section{摸鱼儿}
\textbf{辛弃疾}\par\vspace{0.5em}
更能消、几番风雨。\par
匆匆春又归去。\par
惜春长恨花开早，何况落红无数。\par
春且住。\par
见说道、天涯芳草迷归路。\par
怨春不语。\par
算只有殷勤，画檐珠网，尽日惹飞絮。\par
长门事，准拟佳期又误。\par
蛾眉曾有人妒。\par
千金纵买相如赋，脉脉此情谁诉。\par
君莫舞。\par
君不见、玉环飞燕皆尘土。\par
闲愁最苦。\par
休去倚危楼，斜阳正在，烟柳断肠处。\par

\section{水龙吟}
\textbf{辛弃疾}\par\vspace{0.5em}
楚天千里清秋，水随天去秋无际。\par
遥岑远目，献愁供恨，玉簪螺髻。\par
落日楼头，断鸿声里，江南游子。\par
把吴钩看了，栏干拍遍，无人会、登临意。\par
休说鲈鱼堪。\par
尽西风、季鹰归未。\par
求田问舍，怕应羞见，刘郎才气。\par
可惜流年，忧愁风雨，树犹如此。\par
倩何人，唤取盈盈翠袖，英雄泪。\par

\section{念奴娇}
\textbf{辛弃疾}\par\vspace{0.5em}
野棠花落，又匆匆、过了青明时节。\par
地东风欺客梦，一夜云屏寒怯。\par
曲岸持觞，垂杨系马，此地曾轻别。\par
楼空人去，旧游飞燕能说。\par
闻道绮陌东头，行人长见，帘底纤纤月。\par
旧恨春江流未断，新恨云山千叠。\par
料得明朝，尊前重见，镜里花难折。\par
也应惊问，近来多少华发。\par

\section{鹧鸪天}
\textbf{辛弃疾}\par\vspace{0.5em}
枕簟溪堂冷欲秋。\par
断云依水晚来收。\par
红莲相倚浑如醉，白鸟无言定自愁。\par
书咄咄，且休休。\par
一丘一壑也风流。\par
不知筋力衰多少，但觉新来懒上楼。\par

\section{菩萨蛮}
\textbf{辛弃疾}\par\vspace{0.5em}
郁孤台下清江水。\par
中间多少行人泪。\par
西北是长安。\par
可怜无数山。\par
青山遮不住。\par
毕竟江流去。\par
江晚正愁予。\par
山深闻鹧鸪。\par

\section{木兰花慢}
\textbf{辛弃疾}\par\vspace{0.5em}
老来情味减，对别酒、怯流年。\par
况屈指中秋，十分好月，不照人圆。\par
无情水、都不管，共西风、只等送归船。\par
秋晚莼鲈江上，夜深儿女灯前。\par
征衫。\par
便好去朝天。\par
玉殿正思贤。\par
想夜半承明，留教视草，却遣筹边。\par
长安故人问我，道寻常、泥酒只依然。\par
目断秋霄落雁，醉来时响空弦。\par

\section{祝英台令}
\textbf{辛弃疾}\par\vspace{0.5em}
宝钗分，桃叶渡。\par
烟柳暗南浦。\par
怕上层楼，十日九风雨。\par
断肠片片飞红，都无人管，倩谁唤、流莺声住。\par
鬓边觑。\par
试把花卜心期，才簪又重数。\par
罗帐灯昏，呜咽梦中语。\par
是他春带愁来，春归何处。\par
却不解、将愁归去。\par

\section{青玉案}
\textbf{辛弃疾}\par\vspace{0.5em}
东风夜放花千树。\par
更吹落、星如雨。\par
宝马雕车香满路。\par
凤箫声动，玉壶光转，一夜鱼龙舞。\par
蛾儿雪柳黄金缕。\par
笑语盈盈暗香去。\par
众里寻他千百度。\par
蓦然回首，那人却在，灯火阑珊处。\par

\section{贺新郎}
\textbf{辛弃疾}\par\vspace{0.5em}
凤尾龙香拨。\par
自开元、霓裳曲罢，几番风月。\par
最苦浔阳江头客，画舸亭亭待发。\par
记出塞、黄云堆雪。\par
马上离愁三万里，望昭阳、宫殿孤鸿没。\par
弦解语，恨难说。\par
辽阳驿使音尘绝。\par
琐窗寒、轻拢慢拈，泪珠盈睫。\par
推手含情还却手，一抹梁州哀彻。\par
千古事、云飞烟灭。\par
贺老定场无消息，想沈香亭北繁华歇。\par
弹到此，为呜咽。\par

\section{贺新郎}
\textbf{辛弃疾}\par\vspace{0.5em}
绿树听鹈。\par
更那堪、鹧鸪声住，杜鹃声切。\par
啼到春归无寻处，苦恨芳菲都歇。\par
算未抵、人间离别。\par
马上琵琶关塞黑，更长门、翠辇辞金阙。\par
看燕燕，送归妾。\par
将军百战身名裂。\par
向河梁、回头万里，故人长绝。\par
易水萧萧西风冷，满座衣冠似雪。\par
正壮士、悲歌未彻。\par
啼鸟还知如许恨，料不啼清泪长啼血。\par
谁共我，醉明月。\par

\section{汉宫春}
\textbf{辛弃疾}\par\vspace{0.5em}
春已归来，看美人头上，袅袅春幡。\par
无端风雨，未肯收尽馀寒。\par
年时燕子，料今宵、梦到西园。\par
浑未办、黄柑荐酒，更传青韭堆盘。\par
却笑东风从此，便薰梅染柳，更没些闲。\par
闲时又来镜里，转变朱颜。\par
清愁不断，问何人、会解连环。\par
生怕见、花开花落，朝来塞雁先还。\par

\section{水龙吟}
\textbf{程垓}\par\vspace{0.5em}
夜来风雨匆匆，故园定是花无几。\par
愁多愁极，等闲孤负，一年芳意。\par
柳困花慵，杏青梅小，对人容易。\par
算好春长在，好花长见，元只是、人憔悴。\par
回首池南旧事。\par
恨星星、不堪重记。\par
如今但有，看花老眼，伤时清泪。\par
不怕逢花瘦，只愁怕、老来风味。\par
待繁红乱处，留云借月，也须拼醉。\par

\section{绿头鸭・多丽}
\textbf{贺铸}\par\vspace{0.5em}
玉人家，画楼珠箔临津。\par
托微风、彩箫流怨，断肠马上曾闻。\par
燕堂开、艳妆丛里，调琴思、认歌颦。\par
麝蜡烟浓，玉莲漏短，更衣不待酒初醺。\par
绣屏掩、枕鸳相就，香气渐暾暾。\par
回廊影，疏钟淡月，几许销魂。\par
翠钗分、银笺封泪，舞鞋从此生尘。\par
住兰舟、载将离恨，转南浦、背西曛。\par
记取明年，蔷薇谢後，佳期应未误行云。\par
凤城远，楚梅香嫩，先寄一枝春。\par
青门外，祗凭芳草，寻访郎君。\par

\section{水龙吟}
\textbf{晁补之}\par\vspace{0.5em}
问春何苦匆匆，带风伴雨如驰骤。\par
幽葩细萼，小园低槛，壅培未就。\par
吹尽繁红，占春长久，不如垂柳。\par
算春常不老，人愁春老，愁只是、人间有。\par
春恨十常八九。\par
忍轻辜、芳醪经口。\par
那知自是，桃花结子，不因春瘦。\par
世上功名，老来风味，春归时候。\par
纵樽前痛饮，狂歌似旧，情难依旧。\par

\section{忆少年}
\textbf{晁补之}\par\vspace{0.5em}
无穷官柳，无情画舸，无根行客。\par
南山尚相送，只高城人隔。\par
罨画园林溪绀碧。\par
算重来、尽成陈迹。\par
刘郎鬓如此，况桃花颜色。\par

\section{洞仙歌}
\textbf{晁补之}\par\vspace{0.5em}
青烟幂处，碧海飞金镜。\par
永夜闲阶卧桂影。\par
露凉时、零乱多少寒，神京远，惟有蓝桥路近。\par
水晶帘不下，云母屏开，冷浸佳人淡脂粉。\par
待都将许多明，付与金尊，投晓共、流霞倾尽。\par
更携取、胡床上南楼，看玉做人间，素秋千倾。\par

\section{瑞龙吟}
\textbf{周邦彦}\par\vspace{0.5em}
章台路。\par
还见褪粉梅梢，试花桃树。\par
坊陌人家，定巢燕子，归来旧处。\par
暗凝伫。\par
因念个人痴小，乍窥门户。\par
侵晨浅约宫黄，障风映袖，盈盈笑语。\par
前度刘郎重到，访邻寻里，同时歌舞。\par
唯有旧家秋娘，声价如故。\par
吟笺赋笔，犹记燕台句。\par
知谁伴，名园露饮，东城闲步。\par
事与孤鸿去。\par
探春尽是，伤离意绪。\par
官柳低金缕。\par
归骑晚，纤纤池塘飞雨。\par
断肠院落，一帘风絮。\par

\section{锁窗寒・琐寒窗}
\textbf{周邦彦}\par\vspace{0.5em}
暗柳啼鸦，单衣伫立，小帘朱户。\par
桐花半亩，静锁一庭愁雨。\par
洒空阶，夜阑未休，故人剪烛西窗语。\par
似楚江暝宿，风灯零乱，少年羁旅。\par
迟暮。\par
嬉游处。\par
正店舍无烟，禁城百五。\par
旗亭唤酒，付与高阳俦侣。\par
想东园，桃李自春，小唇秀靥今在否。\par
到归时，定有残英，待客携尊俎。\par

\section{风流子}
\textbf{周邦彦}\par\vspace{0.5em}
新绿小池塘。\par
风帘动，碎影舞斜阳。\par
羡金屋去来，旧时巢燕，土花缭绕，前度莓墙。\par
绣阁凤帏深几许，曾听得理丝簧。\par
欲说又休，虑乖芳信，未歌先咽，愁近清觞。\par
遥知新妆了，开朱户，应自待月西厢。\par
最苦梦魂，今宵不到伊行。\par
问甚时说与，佳音密耗，寄将秦镜，偷换韩香。\par
天便教人，霎时厮见何妨。\par

\section{应天长}
\textbf{周邦彦}\par\vspace{0.5em}
条风布暖，霏雾弄晴，池塘遍满春色。\par
正是夜堂无月，沈沈暗寒食。\par
梁间燕，前社客。\par
似笑我、闭门愁寂。\par
乱花过，隔院芸香，满地狼藉。\par
长记那回时，邂逅相逢，郊外驻油壁。\par
又见汉宫传烛，飞烟五候宅。\par
青青草，迷路陌。\par
强带酒，细寻前迹。\par
市桥远，柳下人家，犹自相识。\par

\section{解连环}
\textbf{周邦彦}\par\vspace{0.5em}
怨怀无托。\par
嗟情人断绝，信音辽邈。\par
信妙手、能解连环，似风散雨收，雾轻云薄。\par
燕子楼空，暗尘锁、一床弦索。\par
想移根换叶。\par
尽是旧时，手种红药。\par
汀洲渐生杜若。\par
料舟依岸曲，人在天角。\par
谩记得、当日音书，把闲语闲言，待总烧却。\par
水驿春回，望寄我、江南梅萼。\par
拚今生，对花对酒，为伊落泪。\par

\section{瑞鹤仙}
\textbf{周邦彦}\par\vspace{0.5em}
悄郊原带郭。\par
行路永，客去车尘漠漠。\par
斜阳映山落。\par
敛馀红、犹恋孤城栏角。\par
凌波步弱。\par
过短亭、何用素约。\par
有流莺劝我，重解绣鞍，缓引春酌。\par
不记归时早暮，上马谁扶，醒眠朱阁。\par
惊飙动幕。\par
扶残醉，绕红药。\par
叹西园、已是花深无地，东风何事又恶。\par
任流光过却。\par
犹喜洞天自乐。\par

\section{浪涛沙・浪淘沙}
\textbf{周邦彦}\par\vspace{0.5em}
昼阴重，霜凋岸草，雾隐城堞。\par
南陌脂车待发。\par
东门帐饮乍阕。\par
正拂面垂杨堪缆结。\par
掩红泪、玉手亲折。\par
念汉浦离鸿去何许，经时信音绝。\par
情切。\par
望中地远天阔。\par
向露冷风清，无人处、耿耿寒漏咽。\par
嗟万事难忘，唯是轻别。\par
翠尊未竭。\par
凭断云留取，西楼残月。\par
罗带光销纹衾叠。\par
连环解、旧香顿歇。\par
怨歌永、琼壶敲尽缺。\par
恨春去、不与人期，弄夜色，空馀满地梨花雪。\par

\section{满庭芳}
\textbf{周邦彦}\par\vspace{0.5em}
风老莺雏，雨肥梅子，午阴嘉树清圆。\par
地卑山近，衣润费垆烟。\par
人静乌鸢自乐，小桥外、新绿溅溅。\par
凭栏久，黄芦苦竹，拟泛九江船。\par
年年。\par
如社燕，飘流瀚海，来寄修椽。\par
且莫思身外，长近尊前。\par
憔悴江南倦客，不堪听、急管繁弦。\par
歌筵畔，先安簟枕，容我醉时眠。\par

\section{过秦楼}
\textbf{周邦彦}\par\vspace{0.5em}
水浴清蟾，叶喧凉吹，巷陌马声初断。\par
闲依露井，笑扑流萤，惹破画罗轻扇。\par
人静夜久凭阑，愁不归眠，立残更箭。\par
叹年华一瞬，人今千里，梦沈书远。\par
空见说、鬓怯琼梳，容销金镜，渐懒趁时匀染。\par
梅风地溽，虹雨苔滋，一架舞红都变。\par
谁信无，为伊才减江淹，情伤荀倩。\par
但明河影下，还看稀星数点。\par

\section{夜游宫}
\textbf{周邦彦}\par\vspace{0.5em}
叶下斜阳照水。\par
卷轻浪、沈沈千里。\par
桥上酸风射眸子。\par
立多时，看黄昏，灯火市。\par
古屋寒窗底。\par
听几片、井桐飞坠。\par
不恋单衾再三起。\par
有谁知，为萧娘，书一纸。\par

\section{解语花}
\textbf{周邦彦}\par\vspace{0.5em}
风销焰蜡，露烘炉。\par
花市光相射。\par
桂华流瓦。\par
纤云散，耿耿素娥欲下。\par
衣裳淡雅。\par
看楚女、纤腰一把。\par
箫鼓喧，人影参差，满路飘香麝。\par
因念都城放夜。\par
望千门如昼，嬉笑游冶。\par
钿车罗帕。\par
相逢处，自有暗尘随马。\par
年光是也。\par
唯只见、旧情衰谢。\par
清漏移，飞盖归来，从舞休歌罢。\par

\section{大}
\textbf{周邦彦}\par\vspace{0.5em}
对宿烟收，春禽静，飞雨时鸣高屋。\par
墙头青玉旆，洗铅霜都尽，嫩梢相触。\par
润逼琴丝，寒侵枕障，虫网吹沾帘竹。\par
邮亭无人处，听檐声不断，困眠初熟。\par
奈愁极顿惊，梦轻难记，自怜幽独。\par
行人归意速。\par
最先念、流潦妨车毂。\par
怎奈向、兰成憔悴，卫清羸，等闲时、易伤心目。\par
未怪平阳客，双泪落、笛中哀曲。\par
况萧索、青芜国。\par
红糁铺地，门外荆桃如菽。\par
夜游共谁秉烛。\par

\section{花犯}
\textbf{周邦彦}\par\vspace{0.5em}
粉墙低，梅花照眼，依然旧风味。\par
露痕轻缀。\par
疑净洗铅华，无限佳丽。\par
去年胜赏曾孤倚。\par
冰盘同宴喜。\par
更可惜，雪中高树，香篝熏素被。\par
今年对花最匆匆，相逢似有恨，依依愁悴。\par
吟望久，青苔上、旋看飞坠。\par
相将见、脆丸荐酒，人正在、空江烟浪里。\par
但梦想、一枝潇洒，黄昏斜照水。\par

\section{六丑}
\textbf{周邦彦}\par\vspace{0.5em}
正单衣试酒，恨客里、光阴虚掷。\par
愿春暂留，春归如过翼。\par
一去无迹。\par
为问花何在，夜来风雨，葬楚宫倾国。\par
钗钿堕处遗香泽。\par
乱点桃蹊，轻翻柳陌。\par
多情为谁追惜。\par
但蜂媒蝶使，时叩窗隔。\par
东园岑寂。\par
渐蒙笼暗碧。\par
静绕珍丛底，成叹息。\par
长条故惹行客。\par
似牵衣待话，别情无极。\par
残英小、强簪巾帻。\par
终不似一朵，钗头颤袅，向人侧。\par
漂流处、莫趁潮汐。\par
恐断红、尚有相思字，何由见得。\par

\section{兰陵王}
\textbf{周邦彦}\par\vspace{0.5em}
柳阴直。\par
烟里丝丝弄碧。\par
隋堤上、曾见几番，拂水飘绵送行色。\par
登临望故国。\par
谁识。\par
京华倦客。\par
长亭路，年去岁来，应折柔条过千尺。\par
闲寻旧踪迹。\par
又酒趁哀弦，灯照离席。\par
梨花榆火催寒食。\par
愁一箭风快，半篙波暖，回头迢递便数驿。\par
望人在天北。\par
凄恻。\par
恨堆积。\par
渐别浦萦回，津堠岑寂。\par
斜阳冉冉春无极。\par
念月榭携手，露桥闻笛。\par
沈思前事，似梦里，泪暗滴。\par

\section{西河}
\textbf{周邦彦}\par\vspace{0.5em}
佳丽地。\par
南朝盛事谁记。\par
山围故国绕清江，髻鬟对起。\par
怒涛寂寞打孤城，风樯遥度天际。\par
断崖树，犹倒倚。\par
莫愁艇子曾系。\par
空馀旧迹郁苍苍，雾沈半垒。\par
夜深月过女墙来，赏心东望淮水。\par
酒旗戏鼓甚处市。\par
想依稀、王谢邻里。\par
燕子不知何世。\par
入寻常、巷陌人家，相对如说兴亡，斜阳里。\par

\section{绮寮怨}
\textbf{周邦彦}\par\vspace{0.5em}
上马人扶残醉，晓风吹未醒。\par
映水曲、翠瓦朱檐，垂杨里、乍见津亭。\par
当时曾题败壁，蛛丝罩、淡墨苔晕青。\par
念去来、岁月如流，徘徊久、叹息愁思盈。\par
去去倦寻路程。\par
江陵旧事，何曾再问杨琼。\par
旧曲凄清。\par
敛愁黛、与谁听。\par
尊前故人如在，想念我、最关情。\par
何须渭城。\par
歌声未尽处，先泪零。\par

\section{拜星月・拜星月慢}
\textbf{周邦彦}\par\vspace{0.5em}
夜色催更，清尘收露，小曲幽坊月暗。\par
竹槛灯窗，识秋娘庭院。\par
笑相遇，似觉琼枝玉树相倚，暖日明霞光烂。\par
水眄兰情，总平生稀见。\par
画图中、旧识春风面。\par
谁知道、自到瑶台畔。\par
眷恋雨润云温，苦惊风吹散。\par
念荒寒、寄宿无人馆。\par
重门闭、败壁秋虫叹。\par
怎奈向、一缕相思，隔溪山不断。\par

\section{尉迟杯}
\textbf{周邦彦}\par\vspace{0.5em}
隋堤路。\par
渐日晚、密霭生深树。\par
阴阴淡月笼沙，还宿河桥深处。\par
无情画舸，都不管、烟波隔南浦。\par
等行人、醉拥重衾，载将离恨归去。\par
因念旧客京华，长偎傍、疏林小槛欢聚。\par
冶叶倡条俱相识，仍惯见、珠歌翠舞。\par
如今向、渔村水驿，夜如岁、焚香独自语。\par
有何人、念我无，梦魂凝想鸳侣。\par

\section{蝶恋花}
\textbf{周邦彦}\par\vspace{0.5em}
月皎惊乌栖不定。\par
更漏将残，辘，牵金井。\par
唤起两眸清炯炯。\par
泪花落枕红棉泠。\par
执手霜风吹鬓影。\par
去意徊徨，别语愁难听。\par
楼上阑干横斗柄。\par
露寒人远鸡相应。\par

\section{夜飞鹊・夜飞鹊慢}
\textbf{周邦彦}\par\vspace{0.5em}
河桥送人处，凉夜何其。\par
斜月远堕馀辉。\par
铜盘烛泪已流尽，霏霏凉露沾衣。\par
相将散离会，探风前津鼓，树杪参旗。\par
华骢会意，纵扬鞭、亦自行迟。\par
迢递路回清野，人语渐无闻，空带愁归。\par
何意重红满地，遗钿不见，斜径都迷。\par
兔葵燕麦，向残阳、欲与人齐。\par
但徘徊班草，欷酹酒，极望天西。\par

\section{关河令・清商怨}
\textbf{周邦彦}\par\vspace{0.5em}
秋阴时晴渐向暝，变一庭凄冷。\par
伫听寒声，云深无雁影。\par
更深人去寂静。\par
但照壁、孤灯相映。\par
酒已都醒，如何消夜永。\par

\section{南浦}
\textbf{孔夷}\par\vspace{0.5em}
风悲画角，听单于、三弄落谯门。\par
投宿征骑，飞雪满孤村。\par
酒市渐闲灯火，正敲窗、乱叶舞纷纷。\par
送数声惊雁，下离烟水，嘹唳度寒云。\par
好在半胧溪月，到如今、无处不销魂。\par
故国梅花归梦，愁损绿罗裙。\par
为问暗香闲艳，也相思、万点付啼痕。\par
算翠屏应是，两眉馀恨倚黄昏。\par

\section{临江仙}
\textbf{晁冲之}\par\vspace{0.5em}
忆昔西池池上饮，年年多少欢娱。\par
别来不寄一行书。\par
寻常相见了，犹道不如初。\par
安稳锦屏今夜梦，月明好度江湖。\par
相思休问定何如。\par
情知春去後，管得落花无。\par

\section{惜分飞}
\textbf{毛滂}\par\vspace{0.5em}
泪湿阑干花著露。\par
愁到眉峰碧聚。\par
此恨平分取。\par
更无言语。\par
空相觑。\par
短雨残云无意绪。\par
寂寞朝朝暮暮。\par
今夜山深处。\par
断魂分付。\par
潮回去。\par

\section{天香}
\textbf{王沂孙}\par\vspace{0.5em}
孤峤蟠烟，层涛蜕月，骊宫夜采铅水。\par
讯远槎风，梦深薇露，化作断魂心字。\par
红瓷候火，还乍识、冰环玉指。\par
一缕萦帘翠影，依稀海天云气。\par
几回娇半醉。\par
翦春灯、夜寒花碎。\par
更好故溪飞雪，小窗深闭。\par
荀令如今顿老，总忘却、樽前旧风味。\par
谩惜馀熏，空篝素被。\par

\section{眉妩}
\textbf{王沂孙}\par\vspace{0.5em}
渐新痕悬柳，澹彩穿花，依约破初暝。\par
便有团圆意，深深拜，相逢谁在香径。\par
画眉未稳，料素娥、犹带离恨。\par
最堪爱、一曲银钩小，宝帘挂秋冷。\par
千古盈亏休问。\par
叹慢磨玉斧，难补金镜。\par
太液池犹在，凄凉处、何人重赋清景。\par
故山夜永。\par
试待他、窥户端正。\par
看云外山河，还老尽、桂花影。\par

\section{齐天乐}
\textbf{王沂孙}\par\vspace{0.5em}
一襟馀恨宫魂断，年年翠阴庭树。\par
乍咽凉柯，还移暗叶，重把离愁深诉。\par
西窗过雨。\par
怪瑶佩流空，玉筝调柱。\par
镜暗妆残，为谁娇鬓尚如许。\par
铜仙铅泪似洗，叹携盘去远，难贮零露。\par
病翼惊秋，枯形阅世，消得斜阳几度。\par
馀音更苦。\par
甚独抱清高，顿成凄楚。\par
谩想薰风，柳丝千万缕。\par

\section{高阳台}
\textbf{王沂孙}\par\vspace{0.5em}
残雪庭阴，轻寒帘影，霏霏玉管春葭。\par
小帖金泥，不知春在谁家。\par
相思一夜窗前梦，奈个人、水隔天遮。\par
但凄然，满树幽香，满地横斜。\par
江南自是离愁苦，况游骢古道，归雁平沙。\par
怎得银笺，殷勤与说年华。\par
如今处处生芳草，纵凭高、不见天涯。\par
更消他，几度东风，几度飞花。\par

\section{法曲献仙音}
\textbf{王沂孙}\par\vspace{0.5em}
层绿峨峨，纤琼皎皎，倒压波痕清浅。\par
过眼年华，动人幽意，相逢几番春换。\par
记唤酒寻芳处，盈盈褪妆晚。\par
已销黯。\par
况凄凉、近来离思，应忘却、明月夜深归辇。\par
荏苒一枝春，恨东风、人似天远。\par
纵有残花，洒征衣、铅泪都满。\par
但殷勤折取，自遣一襟幽怨。\par

\section{长亭怨・长亭怨慢}
\textbf{王沂孙}\par\vspace{0.5em}
泛孤艇、东皋过遍。\par
尚记当日，绿阴门掩。\par
屐齿莓阶，酒痕罗袖事何限。\par
欲寻前迹，空惆怅、成秋苑。\par
自约赏花人，别後总、风流云散。\par
水远。\par
怎知流水外，却是乱山尤远。\par
天涯梦短。\par
想忘了、绮疏雕槛。\par
望不尽、苒苒斜阳，抚乔木、年华将晚。\par
但数点红英，犹识西园凄婉。\par

\section{紫萸香慢}
\textbf{姚云文}\par\vspace{0.5em}
近重阳、偏多风雨，绝怜此日暄明。\par
问秋香浓未，待携客、出西城。\par
正自羁怀多感，怕荒台高处，更不胜情。\par
向尊前、又忆洒酒插花人。\par
只座上、已无老兵。\par
凄情。\par
浅醉还醒。\par
愁不肯、与诗平。\par
记长楸走马，雕弓笮柳，前事休评。\par
紫萸一枝传赐，梦谁到、汉家陵。\par
尽乌纱、便随风去，要天知道，华发如此星星。\par
歌罢涕零。\par

\section{贺新郎}
\textbf{蒋捷}\par\vspace{0.5em}
梦冷黄金屋。\par
叹秦筝、斜鸿阵里，素弦尘扑。\par
化作娇莺飞归去，犹认纱窗旧绿。\par
正过雨、荆桃如菽。\par
此恨难平君知否，似琼台、涌起弹棋局。\par
消瘦影，嫌明烛。\par
鸳楼碎泻东西玉。\par
问芳、何时再展，翠钗难卜。\par
待把宫眉横云样，描上生绡画幅。\par
怕不是、新来妆束。\par
彩扇红牙今都在，恨无人、解听开元曲。\par
空掩袖，倚寒竹。\par

\section{女冠子}
\textbf{蒋捷}\par\vspace{0.5em}
蕙花香也。\par
雪晴池馆如画。\par
春风飞到，宝钗楼上，一片笙箫，琉璃光射。\par
而今灯漫挂。\par
不是暗尘明月，那时元夜。\par
况年来、心懒意怯，羞与蛾儿争要。\par
江城人悄初更打。\par
问繁华谁解，再向天公借。\par
剔残红灺。\par
但梦里隐隐，钿车罗帕。\par
吴笺银粉砑。\par
待把旧家风景，写成闲话。\par
笑绿鬟邻女，倚窗犹唱，夕阳西下。\par

\section{瑞鹤仙}
\textbf{蒋捷}\par\vspace{0.5em}
绀烟迷雁迹。\par
渐断鼓零钟，街喧初息。\par
风檠背寒壁。\par
放冰蜍飞到，丝丝窗隙。\par
琼瑰暗泣。\par
念乡关、霜芜似织。\par
漫将身、化鹤归来，忘却旧游端的。\par
欢极。\par
蓬壶蕖浸，花院梨溶，醉连春夕。\par
柯云罢弈。\par
樱桃在，梦难觅。\par
劝清光，乍可幽窗相伴，休照红楼夜笛。\par
怕人间、换谱伊凉，素娥未识。\par

\section{甘州・八声甘州}
\textbf{张炎}\par\vspace{0.5em}
记玉关、踏雪事清游。\par
寒气脆貂裘。\par
傍枯林古道，长河饮马，此意悠悠。\par
短梦依然江表，老泪洒西州。\par
一字无题处，落叶都愁。\par
载取白云归去，问谁留楚佩，弄影中洲。\par
折芦花赠远，零落一身秋。\par
向寻常野桥流水，待招来、不是旧沙鸥。\par
空怀感，有斜阳处，却怕登楼。\par

\section{渡江云}
\textbf{张炎}\par\vspace{0.5em}
山空天入海，倚楼望极，风急暮潮初。\par
一帘鸠外雨，几处闲田，隔水动春锄。\par
新烟禁柳，想如今、绿到西湖。\par
犹记得、当年深隐，门掩两三株。\par
愁余。\par
荒洲古溆，断梗疏萍，更漂流何处。\par
空自觉、围羞带减，影怯灯孤。\par
常疑即见桃花面，甚近来、翻笑无书。\par
书纵远，如何也都无。\par

\section{解连环}
\textbf{张炎}\par\vspace{0.5em}
楚江空晚。\par
怅离群万里，恍然惊散。\par
自顾影、欲下寒塘，正沙净草枯，水平天远。\par
写不成书，只寄得、相思一点。\par
料因循误了，残毡拥雪，故人心眼。\par
谁怜旅愁荏苒。\par
谩长门夜悄，锦筝弹怨。\par
想伴侣、犹宿芦花，也曾念春前，去程应转。\par
暮雨相呼，怕蓦地、玉关重见。\par
未羞他、双燕归来，画帘半卷。\par

\section{月下笛}
\textbf{张炎}\par\vspace{0.5em}
万里孤云，清游渐远，故人何处。\par
寒窗梦里，犹记经行旧时路。\par
连昌约略无多柳，第一是、难听夜雨。\par
谩惊回凄悄，相看烛影，拥衾谁语。\par
张绪。\par
归何暮。\par
半零落，依依断桥鸥鹭。\par
天涯倦旅。\par
此时心事良苦。\par
只愁重洒西州泪，问杜曲、人家在否。\par
恐翠袖、正天寒，犹倚梅花那树。\par

\section{点绛唇}
\textbf{苏过}\par\vspace{0.5em}
新月娟娟，夜寒江静山衔斗。\par
起来搔首。\par
梅影横窗瘦。\par
好个霜天，闲却传杯手。\par
君知否。\par
乱鸦啼後。\par
归兴浓如酒。\par

\section{贺新郎}
\textbf{叶梦得}\par\vspace{0.5em}
睡起啼莺语。\par
掩青苔、房栊向晚，乱红无数。\par
吹尽残花无人见，惟有垂杨自舞。\par
渐暖霭、初回轻暑。\par
宝扇重寻明月影，暗尘侵、尚有乘鸾女。\par
惊旧恨，遽如许。\par
江南梦断横江渚。\par
浪黏天、葡萄涨绿，半空烟雨。\par
无限楼前沧波意，谁采苹花寄取。\par
但怅望、兰舟容与。\par
万里云帆何时到，送孤鸿、目断千山阻。\par
谁为我，唱金缕。\par

\section{虞美人}
\textbf{叶梦得}\par\vspace{0.5em}
落花已作风前舞。\par
又送黄昏雨。\par
晓来庭院半残红。\par
惟有游丝千丈、晴空。\par
殷勤花下同携手。\par
更尽杯中酒。\par
美人不用敛蛾眉。\par
我亦多情、无奈酒阑时。\par

\section{喜迁莺}
\textbf{刘一止}\par\vspace{0.5em}
晓光催角。\par
听宿鸟未惊，邻鸡先觉。\par
迤逦烟村，马嘶人起，残月尚穿林薄。\par
泪痕带霜微凝，酒力冲寒犹弱。\par
叹倦客、悄不禁，重染风尘京洛。\par
追念，人别後，心事万重，难觅孤鸿托。\par
翠幌娇深，曲屏香暖，争念岁寒飘泊。\par
怨月恨花烦恼，不是不曾经著。\par
这情味，望一成消减，新来还恶。\par

\section{点绛唇}
\textbf{汪藻}\par\vspace{0.5em}
新月娟娟，夜寒江静山衔斗。\par
起来搔首。\par
梅影横窗瘦。\par
好个霜天，闲却传杯手。\par
君知否。\par
乱鸦啼後。\par
归兴浓於酒。\par

\section{蓦山溪}
\textbf{曹组}\par\vspace{0.5em}
洗妆真态，不作铅华御。\par
竹外一枝斜，想佳人、天寒日暮。\par
黄昏小院，无处著清香，风细细，雪垂垂，何况江头路。\par
月边疏影，梦到消魂处。\par
结子欲黄时，又须著、廉纤细雨。\par
孤芳一世，供断有情愁，销瘦却，东阳也，试问花知否。\par

\section{三台}
\textbf{万俟咏}\par\vspace{0.5em}
见梨花初带夜月，海棠半含朝雨。\par
内苑春、不禁过青门，御沟涨、潜通南浦。\par
东风静、细柳垂金缕。\par
望凤阙、非烟非雾。\par
好时代、朝野多欢，遍九陌、太平箫鼓。\par
乍莺儿百啭断续，燕子飞来飞去。\par
近绿水、台榭映秋千，斗草聚、双双游女。\par
饧香更、酒冷踏青路。\par
会暗识、夭桃朱户。\par
向晚骤、宝马雕鞍，醉襟惹、乱花飞絮。\par
正轻寒轻暖漏永，半阴半晴云暮。\par
禁火天、已是试新妆，岁华到、三分佳处。\par
清明看、汉宫传蜡炬。\par
散翠烟、飞入槐府。\par
敛兵卫、阊阖门开，住传宣、又还休务。\par

\section{江神子慢・江城子慢}
\textbf{田为}\par\vspace{0.5em}
玉台挂秋月。\par
铅素浅，梅花傅香雪。\par
冰姿洁。\par
金莲衬、小小凌波罗袜。\par
雨初歇。\par
楼外孤鸿声渐远，远山外、行人音信绝。\par
此恨对语犹难，那堪更寄书说。\par
教人红销翠减，觉衣宽金缕，都为轻别。\par
太情切。\par
消魂处、画角黄昏时节。\par
声呜咽。\par
落尽庭花春去也，银蟾迥、无情圆又缺。\par
恨伊不似馀香，惹鸳鸯结。\par

\section{菩萨蛮}
\textbf{陈克}\par\vspace{0.5em}
赤阑桥尽香街直。\par
笼街细柳娇无力。\par
金碧上青空。\par
花晴帘影红。\par
黄衫飞白马。\par
日日青楼下。\par
醉眼不逢人。\par
午香吹暗尘。\par

\section{菩萨蛮}
\textbf{陈克}\par\vspace{0.5em}
绿芜墙绕青苔院。\par
中庭日淡芭蕉卷。\par
蝴蝶上阶飞。\par
烘帘自在垂。\par
玉钩双语燕。\par
宝甃杨花转。\par
几处簸钱声。\par
绿窗春睡轻。\par

\section{鹧鸪天}
\textbf{周紫芝}\par\vspace{0.5em}
一点残红欲尽时。\par
乍凉秋气满屏帏。\par
梧桐叶上三更雨，叶叶声声是别离。\par
调宝瑟，拨金猊。\par
那时同唱鹧鸪词。\par
如今风雨西楼夜，不听清歌也泪垂。\par

\section{踏莎行}
\textbf{周紫芝}\par\vspace{0.5em}
情似游丝，人如飞絮。\par
泪珠阁定空相觑。\par
一溪烟柳万丝垂，无因系得兰舟住。\par
雁过斜阳，草迷烟渚。\par
如今已是愁无数。\par
明朝且做莫思量，如何过得今宵去。\par

\section{青玉案}
\textbf{无名氏}\par\vspace{0.5em}
年年社日停针线。\par
怎忍见、双飞燕。\par
今日江城春已半。\par
一身犹在，乱山深处，寂寞溪桥畔。\par
春衫著破谁针线。\par
点点行行泪痕满。\par
落日解鞍芳草岸。\par
花无人戴，酒无人劝，醉也无人管。\par

\section{秦楼月・忆秦娥}
\textbf{范成大}\par\vspace{0.5em}
楼阴缺。\par
阑干影卧东厢月。\par
东厢月。\par
一天风露，杏花如雪。\par
隔烟催漏金虬咽。\par
罗帏暗淡灯花结。\par
灯花结。\par
片时春梦，江南天阔。\par

\section{霜天晓角}
\textbf{范成大}\par\vspace{0.5em}
晚晴风歇。\par
一夜春折威。\par
脉脉花疏天淡，云来去、数枝雪。\par
胜绝。\par
愁亦绝。\par
此情谁共说。\par
惟有两行低雁，知人倚、画楼月。\par

\section{眼儿媚}
\textbf{范成大}\par\vspace{0.5em}
酣酣日脚紫烟浮。\par
妍暖破轻裘。\par
困人天色，醉人花气，午梦扶头。\par
春慵恰似春塘水，一片縠纹愁。\par
溶溶泄泄，东风无力，欲皱还休。\par

\section{六州歌头}
\textbf{张孝祥}\par\vspace{0.5em}
长怀望断，关塞莽然平。\par
征尘暗，霜风劲，悄边声。\par
黯销凝。\par
追想当年事，殆天数，非人力，洙泗上，弦歌地，亦膻腥。\par
隔水毡乡，落日牛羊下，区脱纵横。\par
看名王宵猎，骑火一川明。\par
笳鼓悲鸣。\par
遣人惊。\par
念腰间箭，匣中剑，空埃蠹，竟何成。\par
时易失，心徒壮，岁将零。\par
渺神京。\par
干羽方怀远，静烽燧，且休兵。\par
冠盖使，纷驰骛，若为情。\par
闻道中原遗老，常南望，羽葆霓旌。\par
使行人到此，忠愤气填膺。\par
有泪如倾。\par

\section{念奴娇}
\textbf{张孝祥}\par\vspace{0.5em}
洞庭青草，近中秋、更无一点风色。\par
玉鉴琼田三万顷，著我扁舟一叶。\par
素月分辉，明河共影，表里俱澄澈。\par
悠然心会，妙处难与君说。\par
应念岭海经年，孤光自照，肝肺皆冰雪。\par
短发萧骚襟袖冷，稳泛沧浪空阔。\par
尽吸西江，细斟北斗，万象为宾客。\par
扣舷独笑，不知今夕何夕。\par

\end{document}